% Options for packages loaded elsewhere
% Options for packages loaded elsewhere
\PassOptionsToPackage{unicode}{hyperref}
\PassOptionsToPackage{hyphens}{url}
\PassOptionsToPackage{dvipsnames,svgnames,x11names}{xcolor}
%
\documentclass[
  11pt,
]{beaner}
\usepackage{xcolor}
\usepackage[top=20mm,left=20mm,right=20mm,bottom=20mm]{geometry}
\usepackage{amsmath,amssymb}
\setcounter{secnumdepth}{-\maxdimen} % remove section numbering
\usepackage{iftex}
\ifPDFTeX
  \usepackage[T1]{fontenc}
  \usepackage[utf8]{inputenc}
  \usepackage{textcomp} % provide euro and other symbols
\else % if luatex or xetex
  \usepackage{unicode-math} % this also loads fontspec
  \defaultfontfeatures{Scale=MatchLowercase}
  \defaultfontfeatures[\rmfamily]{Ligatures=TeX,Scale=1}
\fi
\usepackage{lmodern}
\ifPDFTeX\else
  % xetex/luatex font selection
  \setmainfont[]{Arial}
\fi
% Use upquote if available, for straight quotes in verbatim environments
\IfFileExists{upquote.sty}{\usepackage{upquote}}{}
\IfFileExists{microtype.sty}{% use microtype if available
  \usepackage[]{microtype}
  \UseMicrotypeSet[protrusion]{basicmath} % disable protrusion for tt fonts
}{}
\makeatletter
\@ifundefined{KOMAClassName}{% if non-KOMA class
  \IfFileExists{parskip.sty}{%
    \usepackage{parskip}
  }{% else
    \setlength{\parindent}{0pt}
    \setlength{\parskip}{6pt plus 2pt minus 1pt}}
}{% if KOMA class
  \KOMAoptions{parskip=half}}
\makeatother
% Make \paragraph and \subparagraph free-standing
\makeatletter
\ifx\paragraph\undefined\else
  \let\oldparagraph\paragraph
  \renewcommand{\paragraph}{
    \@ifstar
      \xxxParagraphStar
      \xxxParagraphNoStar
  }
  \newcommand{\xxxParagraphStar}[1]{\oldparagraph*{#1}\mbox{}}
  \newcommand{\xxxParagraphNoStar}[1]{\oldparagraph{#1}\mbox{}}
\fi
\ifx\subparagraph\undefined\else
  \let\oldsubparagraph\subparagraph
  \renewcommand{\subparagraph}{
    \@ifstar
      \xxxSubParagraphStar
      \xxxSubParagraphNoStar
  }
  \newcommand{\xxxSubParagraphStar}[1]{\oldsubparagraph*{#1}\mbox{}}
  \newcommand{\xxxSubParagraphNoStar}[1]{\oldsubparagraph{#1}\mbox{}}
\fi
\makeatother


\usepackage{longtable,booktabs,array}
\usepackage{calc} % for calculating minipage widths
% Correct order of tables after \paragraph or \subparagraph
\usepackage{etoolbox}
\makeatletter
\patchcmd\longtable{\par}{\if@noskipsec\mbox{}\fi\par}{}{}
\makeatother
% Allow footnotes in longtable head/foot
\IfFileExists{footnotehyper.sty}{\usepackage{footnotehyper}}{\usepackage{footnote}}
\makesavenoteenv{longtable}
\usepackage{graphicx}
\makeatletter
\newsavebox\pandoc@box
\newcommand*\pandocbounded[1]{% scales image to fit in text height/width
  \sbox\pandoc@box{#1}%
  \Gscale@div\@tempa{\textheight}{\dimexpr\ht\pandoc@box+\dp\pandoc@box\relax}%
  \Gscale@div\@tempb{\linewidth}{\wd\pandoc@box}%
  \ifdim\@tempb\p@<\@tempa\p@\let\@tempa\@tempb\fi% select the smaller of both
  \ifdim\@tempa\p@<\p@\scalebox{\@tempa}{\usebox\pandoc@box}%
  \else\usebox{\pandoc@box}%
  \fi%
}
% Set default figure placement to htbp
\def\fps@figure{htbp}
\makeatother





\setlength{\emergencystretch}{3em} % prevent overfull lines

\providecommand{\tightlist}{%
  \setlength{\itemsep}{0pt}\setlength{\parskip}{0pt}}



 


\makeatletter
\@ifpackageloaded{tcolorbox}{}{\usepackage[skins,breakable]{tcolorbox}}
\@ifpackageloaded{fontawesome5}{}{\usepackage{fontawesome5}}
\definecolor{quarto-callout-color}{HTML}{909090}
\definecolor{quarto-callout-note-color}{HTML}{0758E5}
\definecolor{quarto-callout-important-color}{HTML}{CC1914}
\definecolor{quarto-callout-warning-color}{HTML}{EB9113}
\definecolor{quarto-callout-tip-color}{HTML}{00A047}
\definecolor{quarto-callout-caution-color}{HTML}{FC5300}
\definecolor{quarto-callout-color-frame}{HTML}{acacac}
\definecolor{quarto-callout-note-color-frame}{HTML}{4582ec}
\definecolor{quarto-callout-important-color-frame}{HTML}{d9534f}
\definecolor{quarto-callout-warning-color-frame}{HTML}{f0ad4e}
\definecolor{quarto-callout-tip-color-frame}{HTML}{02b875}
\definecolor{quarto-callout-caution-color-frame}{HTML}{fd7e14}
\makeatother
\makeatletter
\@ifpackageloaded{caption}{}{\usepackage{caption}}
\AtBeginDocument{%
\ifdefined\contentsname
  \renewcommand*\contentsname{Table of contents}
\else
  \newcommand\contentsname{Table of contents}
\fi
\ifdefined\listfigurename
  \renewcommand*\listfigurename{List of Figures}
\else
  \newcommand\listfigurename{List of Figures}
\fi
\ifdefined\listtablename
  \renewcommand*\listtablename{List of Tables}
\else
  \newcommand\listtablename{List of Tables}
\fi
\ifdefined\figurename
  \renewcommand*\figurename{Figure}
\else
  \newcommand\figurename{Figure}
\fi
\ifdefined\tablename
  \renewcommand*\tablename{Table}
\else
  \newcommand\tablename{Table}
\fi
}
\@ifpackageloaded{float}{}{\usepackage{float}}
\floatstyle{ruled}
\@ifundefined{c@chapter}{\newfloat{codelisting}{h}{lop}}{\newfloat{codelisting}{h}{lop}[chapter]}
\floatname{codelisting}{Listing}
\newcommand*\listoflistings{\listof{codelisting}{List of Listings}}
\makeatother
\makeatletter
\makeatother
\makeatletter
\@ifpackageloaded{caption}{}{\usepackage{caption}}
\@ifpackageloaded{subcaption}{}{\usepackage{subcaption}}
\makeatother
\usepackage{bookmark}
\IfFileExists{xurl.sty}{\usepackage{xurl}}{} % add URL line breaks if available
\urlstyle{same}
\hypersetup{
  pdftitle={Redacción de Proyectos de Investigación Científica},
  pdfauthor={Edison Achalma},
  colorlinks=true,
  linkcolor={blue},
  filecolor={Maroon},
  citecolor={Blue},
  urlcolor={Blue},
  pdfcreator={LaTeX via pandoc}}


\title{Redacción de Proyectos de Investigación Científica}
\usepackage{etoolbox}
\makeatletter
\providecommand{\subtitle}[1]{% add subtitle to \maketitle
  \apptocmd{\@title}{\par {\large #1 \par}}{}{}
}
\makeatother
\subtitle{Corporacion Académica Universitaria}
\author{Edison Achalma}
\date{2026-01-26}
\begin{document}
\maketitle


\section{GENERACIÓN DE ESTRUCTURA DE MONOGRAFÍA CON
IA}\label{generaciuxf3n-de-estructura-de-monografuxeda-con-ia}

\subsection{Prompt Estructurado para Crear
Capítulos}\label{prompt-estructurado-para-crear-capuxedtulos}

\begin{center}\rule{0.5\linewidth}{0.5pt}\end{center}

\subsection{Prompt para Generar Estructura de
Monografía}\label{prompt-para-generar-estructura-de-monografuxeda}

\textbf{Leyenda de colores:} {Azul = Rol} {Verde = Contexto} {Naranja =
Tarea} {Morado = Formato} {Rojo = Restricciones}

{ROL:} Eres un asesor académico especializado en metodología de
investigación y estructuración de monografías universitarias. Diseñas
estructuras coherentes, lógicas y académicamente rigurosas.

{CONTEXTO:} Monografía de nivel pregrado evaluada según criterios
académicos estándar. Debe facilitar argumentación clara, desarrollo
lógico y presentación sistemática.

{TAREA:} Genera estructura completa (capítulos, subcapítulos,
sub-subcapítulos) para:

\textbf{Tema:} {[}INSERTAR AQUÍ{]}\\
\textbf{Tipo:} {[}Compilación/Investigación{]}\\
\textbf{Extensión:} {[}XX páginas{]}

{FORMATO:} Introducción → Cap. 1-3 (con 1.1, 1.1.1, etc.) → Conclusiones
→ Referencias

{RESTRICCIONES:} 3-4 capítulos \textbar{} 3-5 subcapítulos/cap
\textbar{} 2-4 sub-subcapítulos \textbar{} Máximo 3 niveles \textbar{}
Títulos específicos, no genéricos \textbar{} Progresión lógica
\textbar{} Equilibrio entre capítulos \textbar{} NO metodología en
compilación

{CALIDAD:} Coherencia temática, progresión lógica, especificidad,
equilibrio, profundidad apropiada

{EJEMPLO:} ``Gestión del agua en Ayacucho'' → Cap. 1: Sistemas
tradicionales (1.1 Organización social: 1.1.1 Autoridades, 1.1.2 Normas,
1.1.3 Conflictos \textbar{} 1.2 Conocimientos: 1.2.1 Calendario ritual,
1.2.2 Señas naturales, 1.2.3 Técnicas almacenamiento)

{IMPORTANTE:} Solo estructura (títulos y numeración), NO contenido.

\begin{center}\rule{0.5\linewidth}{0.5pt}\end{center}

\subsection{Uso del Prompt}\label{uso-del-prompt}

\textbf{1.} Copiar prompt completo \textbar{} \textbf{2.} Reemplazar
{[}INSERTAR AQUÍ{]} con tu tema \textbar{} \textbf{3.} Pegar en
ChatGPT/Claude/Gemini \textbar{} \textbf{4.} Revisar y ajustar
\textbar{} \textbf{5.} Validar con asesor

\textbf{Ejemplo de entrada:}

Tema: ``Impacto de redes sociales en salud mental de adolescentes
peruanos'' \textbar{} Tipo: Compilación \textbar{} Extensión: 60 páginas

\textbf{Ejemplo de salida esperada:}

\textbf{Cap. 1: Uso de redes sociales en adolescentes peruanos} → 1.1
Patrones de acceso (1.1.1 Penetración internet, 1.1.2 Plataformas
predominantes, 1.1.3 Tiempo exposición) \textbar{} 1.2 Motivaciones
(1.2.1 Identidad digital, 1.2.2 Socialización, 1.2.3 Validación social)

\begin{center}\rule{0.5\linewidth}{0.5pt}\end{center}

\subsection{Variaciones Específicas}\label{variaciones-especuxedficas}

\textbf{Interdisciplinaria:} Agregar ``Integrar perspectivas de {[}X{]}
y {[}Y{]}''

\textbf{Comparativa:} Agregar ``Comparar {[}A{]} con {[}B{]}
sistemáticamente''

\textbf{Histórica:} Agregar ``Organizar cronológicamente por períodos''

\textbf{Aplicada:} Agregar ``Incluir capítulo final con propuestas
prácticas''

\begin{center}\rule{0.5\linewidth}{0.5pt}\end{center}

\subsection{Prompt para Desarrollar
Capítulos}\label{prompt-para-desarrollar-capuxedtulos}

\textbf{Leyenda:} {Azul = Rol} {Verde = Tarea} {Naranja = Formato} {Rojo
= Restricción}

{ROL:} Experto en {[}TU TEMA{]} que redacta contenido académico
riguroso.

{TAREA:} Desarrolla contenido para Cap. {[}\#{]}: {[}título{]}
\textbar{} Subcap. {[}\#{]}: {[}título{]}

{FORMATO:} {[}X{]} páginas \textbar{} Incluir definiciones, autores
clave, ejemplos \textbar{} Estilo académico formal \textbar{} Indicar
citas como {[}CITA: Autor, año{]}

{RESTRICCIÓN:} NO inventar datos ni autores. Usar {[}AUTOR A{]},
{[}AUTOR B{]} como marcadores.

\begin{center}\rule{0.5\linewidth}{0.5pt}\end{center}

\subsection{Recomendaciones Finales}\label{recomendaciones-finales}

\textbf{HACER}

Revisar la estructura completa antes de empezar a escribir

Validar con tu asesor académico

Usar la IA como asistente, no como autor

Completar las citas con fuentes reales

Mantener coherencia entre capítulos

\textbf{NO HACER}

Copiar directamente el contenido que genera la IA

Aceptar la primera estructura sin revisión crítica

Usar la estructura sin adaptarla a tu tema específico

Olvidar agregar tu análisis y reflexión personal

Incluir más de 3 niveles de profundidad

\begin{center}\rule{0.5\linewidth}{0.5pt}\end{center}

\section{FUNDAMENTOS}\label{fundamentos}

\subsection{Bienvenida y Objetivos}\label{bienvenida-y-objetivos}

\subsubsection{Objetivos de la Clase}\label{objetivos-de-la-clase}

\begin{itemize}
\tightlist
\item
  Dominar los principios de la redacción científica
\item
  Desarrollar habilidades de escritura académica
\item
  Aplicar herramientas de estilo
\end{itemize}

\subsubsection{Indice}\label{indice}

\begin{itemize}
\tightlist
\item
  Fundamentos
\item
  Proceso de Escritura
\item
  Herramientas Prácticas
\end{itemize}

\begin{center}\rule{0.5\linewidth}{0.5pt}\end{center}

\subsection{¿Qué es la Escritura
Científica?}\label{quuxe9-es-la-escritura-cientuxedfica}

\subsubsection{Definición}\label{definiciuxf3n}

\begin{quote}
La escritura científica es una \textbf{práctica especializada} cuyo
objetivo primordial es la \textbf{comunicación eficaz} de ideas y
resultados de investigación a una comunidad.
\end{quote}

\subsubsection{Diferencias}\label{diferencias}

\begin{longtable}[]{@{}ll@{}}
\toprule\noalign{}
Literaria & Científica \\
\midrule\noalign{}
\endhead
\bottomrule\noalign{}
\endlastfoot
Entretener/Expresar & Informar \\
Metáforas/Suspenso & Precisión/Claridad \\
Subjetiva & Objetiva \\
\end{longtable}

\subsubsection{Propósito}\label{propuxf3sito}

\begin{itemize}
\tightlist
\item
  Informar hallazgos
\item
  Transferir conocimiento
\item
  Integrar al saber universal
\item
  Permitir replicación
\end{itemize}

\subsubsection{Impacto}\label{impacto}

La investigación \textbf{NO termina} hasta que se publica en una revista
con revisión por pares.

\begin{center}\rule{0.5\linewidth}{0.5pt}\end{center}

\subsection{Los 4 Pilares de la Redacción
Científica}\label{los-4-pilares-de-la-redacciuxf3n-cientuxedfica}

\subsubsection{1. PRECISIÓN}\label{precisiuxf3n}

\begin{itemize}
\tightlist
\item
  Palabras exactas
\item
  Sin ambigüedades
\item
  Términos técnicos apropiados
\item
  No obligar al lector a ``adivinar''
\end{itemize}

\subsubsection{2. CLARIDAD}\label{claridad}

\begin{itemize}
\tightlist
\item
  Lenguaje sencillo
\item
  Lectura rápida
\item
  Oraciones bien construidas
\item
  Orden lógico
\end{itemize}

\subsubsection{3. BREVEDAD}\label{brevedad}

\begin{itemize}
\tightlist
\item
  Solo información pertinente
\item
  Mínimo de palabras
\item
  Evitar redundancia
\item
  \textbf{Promedio: 18 palabras/oración}
\end{itemize}

\subsubsection{4. FORMALIDAD}\label{formalidad}

\begin{itemize}
\tightlist
\item
  Registro nivelado
\item
  Sin coloquialismos
\item
  Sin localismos
\item
  Lenguaje profesional
\end{itemize}

\begin{center}\rule{0.5\linewidth}{0.5pt}\end{center}

\subsection{Honestidad Intelectual}\label{honestidad-intelectual}

\subsubsection{Principios Éticos}\label{principios-uxe9ticos}

\begin{itemize}
\tightlist
\item
  Decir la verdad\\
\item
  Presentar todos los hechos relevantes\\
\item
  Reconocer contribuciones ajenas\\
\item
  Evitar el \textbf{PLAGIO}\\
\item
  Abierto a la crítica
\end{itemize}

\subsubsection{Tipos de Plagio}\label{tipos-de-plagio}

\begin{itemize}
\tightlist
\item
  \textbf{Literal:} Copiar textualmente sin citar
\item
  \textbf{Adaptado:} Parafrasear sin citar
\item
  \textbf{De ideas:} Usar ideas ajenas sin crédito
\item
  \textbf{Autoplagio:} Reutilizar trabajo propio sin avisar
\end{itemize}

\subsubsection{Consecuencias del Plagio}\label{consecuencias-del-plagio}

\begin{tcolorbox}[enhanced jigsaw, arc=.35mm, rightrule=.15mm, breakable, titlerule=0mm, title=\textcolor{quarto-callout-warning-color}{\faExclamationTriangle}\hspace{0.5em}{Impacto Grave}, toprule=.15mm, left=2mm, colback=white, colbacktitle=quarto-callout-warning-color!10!white, bottomrule=.15mm, colframe=quarto-callout-warning-color-frame, toptitle=1mm, opacityback=0, opacitybacktitle=0.6, coltitle=black, bottomtitle=1mm, leftrule=.75mm]

\begin{itemize}
\tightlist
\item
  Pérdida de credibilidad
\item
  Rechazo de manuscritos
\item
  Sanciones académicas
\item
  Daño a la carrera profesional
\end{itemize}

\end{tcolorbox}

\subsubsection{Cómo Evitarlo}\label{cuxf3mo-evitarlo}

\begin{enumerate}
\def\labelenumi{\arabic{enumi}.}
\tightlist
\item
  Citar \textbf{SIEMPRE} las fuentes
\item
  Usar comillas para citas textuales
\item
  Parafrasear correctamente
\item
  Gestores bibliográficos (Mendeley, Zotero)
\end{enumerate}

\begin{center}\rule{0.5\linewidth}{0.5pt}\end{center}

\subsection{Géneros Académicos: Formación
vs.~Expertos}\label{guxe9neros-acaduxe9micos-formaciuxf3n-vs.-expertos}

\subsubsection{Géneros de Formación}\label{guxe9neros-de-formaciuxf3n}

\emph{(Estudiantes → Evaluación)}

\begin{itemize}
\tightlist
\item
  \textbf{Examen parcial}

  \begin{itemize}
  \tightlist
  \item
    Presencial: tiempo/memoria limitados
  \item
    Domiciliario: género ``bisagra''
  \end{itemize}
\item
  \textbf{Monografía}

  \begin{itemize}
  \tightlist
  \item
    Entrenamiento investigativo
  \item
    Toma de posición teórica
  \item
    Exhaustivo sobre un tema
  \end{itemize}
\item
  \textbf{Informe de lectura}

  \begin{itemize}
  \tightlist
  \item
    Interpretación de textos
  \item
    Sin marco teórico obligatorio
  \end{itemize}
\end{itemize}

\subsubsection{Géneros Expertos}\label{guxe9neros-expertos}

\emph{(Profesionales → Comunicación)}

\begin{itemize}
\tightlist
\item
  \textbf{Artículo científico}

  \begin{itemize}
  \tightlist
  \item
    Primera publicación de resultados
  \item
    Formato IMRyD
  \item
    Revisión por pares
  \end{itemize}
\item
  \textbf{Nota de revisión}

  \begin{itemize}
  \tightlist
  \item
    Meta-análisis
  \item
    Síntesis de literatura
  \end{itemize}
\item
  \textbf{Tesis doctoral}

  \begin{itemize}
  \tightlist
  \item
    Aporte original al conocimiento
  \item
    Investigación exhaustiva
  \end{itemize}
\end{itemize}

\begin{center}\rule{0.5\linewidth}{0.5pt}\end{center}

\subsection{La Monografía}\label{la-monografuxeda}

\subsubsection{Definición}\label{definiciuxf3n-1}

\begin{quote}
Tratado \textbf{especial}, \textbf{exhaustivo} y \textbf{sistemático}
sobre un asunto determinado dentro de una disciplina.
\end{quote}

\subsubsection{Objetivos}\label{objetivos}

\begin{itemize}
\tightlist
\item
  Entrenar capacidad investigativa
\item
  Tomar posición entre posturas diversas
\item
  Manejar marco teórico
\item
  Comunicar conocimientos rigurosamente
\end{itemize}

\subsubsection{Diferencias}\label{diferencias-1}

\begin{itemize}
\tightlist
\item
  \textbf{vs.~Ensayo:} Requiere marco teórico y metodología
\item
  \textbf{vs.~Informe:} Exige posición del autor
\end{itemize}

\subsubsection{Tipos de Monografía}\label{tipos-de-monografuxeda}

\begin{enumerate}
\def\labelenumi{\arabic{enumi}.}
\tightlist
\item
  \textbf{Compilación}

  \begin{itemize}
  \tightlist
  \item
    Analiza y compara autores
  \item
    Opinión personal fundada
  \end{itemize}
\item
  \textbf{Investigación}

  \begin{itemize}
  \tightlist
  \item
    Temas nuevos o poco explorados
  \item
    Aporte original
  \end{itemize}
\item
  \textbf{Experiencia}

  \begin{itemize}
  \tightlist
  \item
    Basada en vivencias profesionales
  \item
    Sustento teórico a la práctica
  \end{itemize}
\end{enumerate}

\subsubsection{Errores Comunes}\label{errores-comunes}

\begin{itemize}
\tightlist
\item
  Tema excesivamente amplio
\item
  Mero ``hilván de fichas''
\item
  Falta de posición personal
\end{itemize}

\begin{center}\rule{0.5\linewidth}{0.5pt}\end{center}

\section{ESTRUCTURA DE MONOGRAFÍA DE
COMPILACIÓN}\label{estructura-de-monografuxeda-de-compilaciuxf3n}

\begin{center}\rule{0.5\linewidth}{0.5pt}\end{center}

\subsection{¿Qué es una Monografía de
Compilación?}\label{quuxe9-es-una-monografuxeda-de-compilaciuxf3n}

\textbf{Definición:}

Trabajo académico que recopila, organiza y sintetiza información de
diversas fuentes sobre un tema específico, sin generar datos originales.

\textbf{Objetivo:}

Presentar un panorama completo y ordenado del estado del conocimiento
sobre un tema, integrando diferentes perspectivas y autores.

\begin{center}\rule{0.5\linewidth}{0.5pt}\end{center}

\subsection{Estructura General: Monografía de
Compilación}\label{estructura-general-monografuxeda-de-compilaciuxf3n}

\textbf{PORTADA}

Datos institucionales

Universidad, facultad, escuela profesional, título, autor, asesor,
lugar, año

\textbf{ÍNDICE}

Navegación

Lista completa de capítulos, subcapítulos y páginas

\textbf{INTRODUCCIÓN}

Presentación

Justificación, objetivos, alcances, limitaciones, metodología de
compilación

\textbf{CAPÍTULOS}

Desarrollo

2-4 capítulos temáticos con subcapítulos (núcleo del trabajo)

\textbf{CONCLUSIONES}

Síntesis

Síntesis integradora, hallazgos principales, reflexiones finales

\textbf{REFERENCIAS}

Fuentes

Bibliografía completa en formato APA, mínimo 15-20 fuentes

\begin{center}\rule{0.5\linewidth}{0.5pt}\end{center}

\subsection{INTRODUCCIÓN: Cómo
Redactar}\label{introducciuxf3n-cuxf3mo-redactar}

\textbf{PÁRRAFO 1: CONTEXTUALIZACIÓN}

{TEMA AMPLIO}

La educación superior en América Latina ha experimentado
transformaciones significativas en las últimas dos décadas. El acceso a
la universidad se ha democratizado, pero persisten desafíos en términos
de calidad educativa y pertinencia de los programas académicos.

{ENFOQUE ESPECÍFICO}

En particular, la formación docente universitaria requiere atención
especial, dado su rol determinante en la calidad de los procesos de
enseñanza-aprendizaje.

\begin{center}\rule{0.5\linewidth}{0.5pt}\end{center}

\subsection{INTRODUCCIÓN: Cómo Redactar
(2)}\label{introducciuxf3n-cuxf3mo-redactar-2}

\textbf{PÁRRAFO 2: JUSTIFICACIÓN}

{IMPORTANCIA}

Esta monografía se justifica por la necesidad de comprender
integralmente los modelos de formación docente universitaria vigentes en
la región.

{VACÍO/NECESIDAD}

Si bien existen estudios sobre aspectos específicos de esta formación,
hace falta una síntesis que integre las diversas perspectivas teóricas y
experiencias documentadas en la literatura especializada.

\begin{center}\rule{0.5\linewidth}{0.5pt}\end{center}

\subsection{INTRODUCCIÓN: Cómo Redactar
(3)}\label{introducciuxf3n-cuxf3mo-redactar-3}

\textbf{PÁRRAFO 3: OBJETIVOS}

{OBJETIVO GENERAL}

El objetivo de esta monografía es analizar y sintetizar las principales
propuestas teóricas y experiencias prácticas sobre formación docente
universitaria en América Latina, período 2010-2025.

{OBJETIVOS ESPECÍFICOS}

Específicamente, se busca: (a) identificar los modelos teóricos
predominantes, (b) caracterizar las experiencias más relevantes, y (c)
establecer tendencias y desafíos emergentes.

\begin{center}\rule{0.5\linewidth}{0.5pt}\end{center}

\subsection{INTRODUCCIÓN: Cómo Redactar
(4)}\label{introducciuxf3n-cuxf3mo-redactar-4}

\textbf{PÁRRAFO 4: METODOLOGÍA Y ESTRUCTURA}

{METODOLOGÍA}

La investigación documental se realizó mediante búsqueda sistemática en
bases de datos académicas (Scopus, SciELO, Redalyc), seleccionando 45
fuentes entre artículos científicos, libros especializados y documentos
institucionales.

{ESTRUCTURA}

El trabajo se organiza en tres capítulos: el primero aborda los
fundamentos teóricos, el segundo analiza experiencias nacionales
destacadas, y el tercero discute tendencias y desafíos actuales.

\begin{center}\rule{0.5\linewidth}{0.5pt}\end{center}

\subsection{CAPÍTULOS: Cómo Redactar el
Desarrollo}\label{capuxedtulos-cuxf3mo-redactar-el-desarrollo}

\textbf{Cada capítulo debe tener:}

\textbf{INTRODUCCIÓN DEL CAPÍTULO}

Párrafo breve que presenta el contenido del capítulo y su conexión con
el anterior

\textbf{SUBCAPÍTULOS (2-4)}

Cada uno desarrolla un aspecto específico del tema del capítulo

\textbf{SÍNTESIS PARCIAL}

Párrafo final que resume los puntos principales del capítulo

\textbf{CONEXIÓN}

Oración que anticipa el siguiente capítulo

\begin{center}\rule{0.5\linewidth}{0.5pt}\end{center}

\subsection{CAPÍTULO: Ejemplo de Redacción
Inicial}\label{capuxedtulo-ejemplo-de-redacciuxf3n-inicial}

\textbf{CAPÍTULO 1: FUNDAMENTOS TEÓRICOS DE LA FORMACIÓN DOCENTE
UNIVERSITARIA}

{PRESENTACIÓN}

Este primer capítulo examina las bases conceptuales que sustentan la
formación docente universitaria.

{CONTENIDO}

Se analizan tres perspectivas teóricas fundamentales: el enfoque de
competencias profesionales, la perspectiva del conocimiento pedagógico
del contenido, y el modelo de desarrollo profesional continuo.

{ORGANIZACIÓN}

Cada enfoque se presenta en un subcapítulo que incluye sus principios
centrales, autores representativos y aplicaciones al contexto
universitario.

\begin{center}\rule{0.5\linewidth}{0.5pt}\end{center}

\subsection{SUBCAPÍTULO: Ejemplo de
Desarrollo}\label{subcapuxedtulo-ejemplo-de-desarrollo}

\textbf{1.1 Enfoque de competencias profesionales docentes}

{DEFINICIÓN}

El enfoque de competencias profesionales docentes surge en la década de
1990 como respuesta a la necesidad de operacionalizar los saberes
requeridos para el ejercicio efectivo de la docencia (Perrenoud, 2004).

{AUTORES CLAVE}

Según Zabalza (2009), las competencias docentes universitarias integran
conocimientos disciplinares, habilidades pedagógicas y actitudes
profesionales que permiten al docente diseñar, desarrollar y evaluar
procesos formativos de calidad.

{COMPONENTES}

Este modelo identifica típicamente entre 8 y 12 competencias clave, que
incluyen el dominio disciplinar, la planificación didáctica, la
evaluación de aprendizajes, y el uso de tecnologías educativas, entre
otras (Mas \& Olmos, 2016).

\begin{center}\rule{0.5\linewidth}{0.5pt}\end{center}

\subsection{CONCLUSIONES: Cómo
Redactar}\label{conclusiones-cuxf3mo-redactar}

\textbf{PÁRRAFO 1: SÍNTESIS INTEGRADORA}

{RECAPITULACIÓN}

Esta monografía ha examinado las principales dimensiones de la formación
docente universitaria en América Latina, integrando aportes teóricos y
experiencias documentadas en la literatura especializada del período
2010-2025.

{HALLAZGO PRINCIPAL}

El análisis realizado evidencia una transición desde modelos centrados
en la transmisión de conocimientos hacia propuestas que enfatizan el
desarrollo de competencias profesionales complejas y el aprendizaje
reflexivo.

\begin{center}\rule{0.5\linewidth}{0.5pt}\end{center}

\subsection{CONCLUSIONES: Cómo Redactar
(2)}\label{conclusiones-cuxf3mo-redactar-2}

\textbf{PÁRRAFO 2: HALLAZGOS POR OBJETIVO}

{OBJETIVO 1}

Respecto a los modelos teóricos, se identificaron tres perspectivas
predominantes que coexisten y se complementan en las propuestas
formativas actuales.

{OBJETIVO 2}

En cuanto a las experiencias prácticas, destacan iniciativas que
combinan formación inicial, acompañamiento in situ y comunidades de
práctica como estrategias efectivas.

{OBJETIVO 3}

Las tendencias emergentes apuntan hacia la internacionalización, la
innovación pedagógica digital y el énfasis en la investigación educativa
como ejes articuladores.

\begin{center}\rule{0.5\linewidth}{0.5pt}\end{center}

\subsection{CONCLUSIONES: Cómo Redactar
(3)}\label{conclusiones-cuxf3mo-redactar-3}

\textbf{PÁRRAFO 3: REFLEXIÓN FINAL}

{IMPLICACIONES}

Los hallazgos de esta compilación sugieren que la calidad de la
formación docente universitaria requiere políticas institucionales
integrales que articulen desarrollo profesional, condiciones laborales
adecuadas y reconocimiento de la función docente.

{PROSPECTIVA}

Futuras investigaciones podrían profundizar en el impacto de estas
propuestas formativas sobre los aprendizajes estudiantiles, así como en
los factores institucionales que facilitan u obstaculizan su
implementación efectiva.

\begin{center}\rule{0.5\linewidth}{0.5pt}\end{center}

\section{ESTRUCTURA DE MONOGRAFÍA DE
INVESTIGACIÓN}\label{estructura-de-monografuxeda-de-investigaciuxf3n}

\begin{center}\rule{0.5\linewidth}{0.5pt}\end{center}

\subsection{¿Qué es una Monografía de
Investigación?}\label{quuxe9-es-una-monografuxeda-de-investigaciuxf3n}

\textbf{Definición:}

Trabajo académico que presenta resultados de una investigación original,
generando datos nuevos mediante trabajo de campo, experimentos o
análisis de fuentes primarias.

\textbf{Diferencia clave:}

No solo recopila información existente, sino que produce conocimiento
nuevo a través de metodología científica aplicada.

\begin{center}\rule{0.5\linewidth}{0.5pt}\end{center}

\subsection{Estructura General: Monografía de
Investigación}\label{estructura-general-monografuxeda-de-investigaciuxf3n}

\textbf{PORTADA}

Identificación

Universidad, facultad, título, autor, asesor, lugar, año

\textbf{RESUMEN}

Síntesis

Abstract ejecutivo (250-300 palabras): problema, método, resultados,
conclusión

\textbf{INTRODUCCIÓN}

Marco general

Problema, justificación, objetivos, hipótesis (si aplica)

\textbf{MARCO TEÓRICO}

Fundamentos

Antecedentes, bases teóricas, definición de conceptos clave

\textbf{METODOLOGÍA}

Diseño

Tipo, enfoque, población, muestra, instrumentos, procedimientos

\textbf{RESULTADOS}

Hallazgos

Presentación organizada de datos obtenidos (tablas, gráficos, análisis)

\textbf{DISCUSIÓN}

Interpretación

Interpretación de resultados, comparación con literatura, limitaciones

\textbf{CONCLUSIONES}

Cierre

Respuesta a objetivos, recomendaciones, futuras investigaciones

\textbf{REFERENCIAS}

Fuentes

Bibliografía completa en formato APA

\textbf{ANEXOS}

Materiales

Instrumentos, transcripciones, datos complementarios

\begin{center}\rule{0.5\linewidth}{0.5pt}\end{center}

\subsection{INTRODUCCIÓN: Cómo Redactar
(Investigación)}\label{introducciuxf3n-cuxf3mo-redactar-investigaciuxf3n}

\textbf{PÁRRAFO 1: PLANTEAMIENTO DEL PROBLEMA}

{CONTEXTO}

Las comunidades campesinas de Ayacucho han desarrollado históricamente
sistemas de gestión del agua basados en conocimientos tradicionales que
regulan el acceso, distribución y conservación de este recurso.

{PROBLEMA}

Sin embargo, estos sistemas enfrentan presiones crecientes debido al
cambio climático, la migración juvenil y la introducción de tecnologías
modernas que no siempre son compatibles con las prácticas locales.

{PREGUNTA}

¿Cómo se articulan los conocimientos tradicionales con las prácticas
contemporáneas de gestión del agua en la comunidad campesina de Chuschi,
Ayacucho?

\begin{center}\rule{0.5\linewidth}{0.5pt}\end{center}

\subsection{INTRODUCCIÓN: Cómo Redactar (Investigación)
(2)}\label{introducciuxf3n-cuxf3mo-redactar-investigaciuxf3n-2}

\textbf{PÁRRAFO 2: OBJETIVOS E HIPÓTESIS}

{OBJETIVO GENERAL}

Analizar los mecanismos de articulación entre conocimientos
tradicionales y prácticas contemporáneas de gestión del agua en Chuschi
durante el período 2023-2024.

{OBJETIVOS ESPECÍFICOS}

\begin{enumerate}
\def\labelenumi{(\alph{enumi})}
\tightlist
\item
  Identificar los conocimientos tradicionales vigentes sobre gestión del
  agua; (b) caracterizar las prácticas contemporáneas introducidas; (c)
  determinar los puntos de tensión y complementariedad entre ambos
  sistemas.
\end{enumerate}

{HIPÓTESIS}

Se plantea que existe un proceso de hibridación selectiva donde la
comunidad integra tecnologías modernas manteniendo los principios
organizativos tradicionales de reciprocidad y gestión colectiva.

\begin{center}\rule{0.5\linewidth}{0.5pt}\end{center}

\subsection{MARCO TEÓRICO: Cómo
Redactar}\label{marco-teuxf3rico-cuxf3mo-redactar}

\textbf{ESTRUCTURA TÍPICA}

{ANTECEDENTES}

Investigaciones previas sobre el tema: ¿Qué se ha estudiado? ¿Qué
hallazgos existen? ¿Qué vacíos persisten?

{BASES TEÓRICAS}

Teorías y enfoques que fundamentan tu investigación: conceptos clave,
modelos explicativos, perspectivas disciplinares.

{DEFINICIONES}

Conceptos operacionales: ¿Cómo defines los términos centrales de tu
investigación?

\begin{center}\rule{0.5\linewidth}{0.5pt}\end{center}

\subsection{MARCO TEÓRICO: Ejemplo de
Redacción}\label{marco-teuxf3rico-ejemplo-de-redacciuxf3n}

\textbf{2.1 Antecedentes de investigación}

{INTERNACIONALES}

Boelens (2008) analizó sistemas de gestión del agua en comunidades
andinas de Ecuador, identificando tensiones entre normativas estatales y
sistemas consuetudinarios locales.

{NACIONALES}

En el contexto peruano, Gelles (2000) documentó prácticas rituales
asociadas a la gestión del agua en Cabanaconde, demostrando la
integración de dimensiones técnicas y simbólicas.

{LOCALES}

Para Ayacucho específicamente, Rénique (2004) estudió cambios en
sistemas agrícolas post-violencia política, aunque sin enfocarse en
gestión hídrica.

{VACÍO}

No se han realizado estudios recientes sobre la articulación
contemporánea de conocimientos tradicionales y modernos en gestión del
agua en comunidades ayacuchanas.

\begin{center}\rule{0.5\linewidth}{0.5pt}\end{center}

\subsection{METODOLOGÍA: Cómo
Redactar}\label{metodologuxeda-cuxf3mo-redactar}

\textbf{PÁRRAFO 1: TIPO Y ENFOQUE}

{TIPO}

La investigación es de tipo básica, nivel descriptivo-explicativo, con
diseño no experimental transversal.

{ENFOQUE}

Se adoptó un enfoque cualitativo con elementos etnográficos,
privilegiando la comprensión profunda de significados y prácticas desde
la perspectiva de los actores locales.

\begin{center}\rule{0.5\linewidth}{0.5pt}\end{center}

\subsection{METODOLOGÍA: Cómo Redactar
(2)}\label{metodologuxeda-cuxf3mo-redactar-2}

\textbf{PÁRRAFO 2: POBLACIÓN Y MUESTRA}

{POBLACIÓN}

La población de estudio estuvo conformada por los 180 comuneros activos
de Chuschi registrados en el padrón comunal 2023.

{MUESTRA}

Se seleccionó una muestra no probabilística intencional de 32 comuneros:
12 autoridades de riego (incluyendo 4 yaku alcaldes), 15 usuarios de
riego con más de 20 años de experiencia, y 5 jóvenes comuneros con
estudios técnicos.

{CRITERIOS}

Los criterios de inclusión fueron: residencia permanente, participación
activa en faenas de agua, y disposición para participar en entrevistas.

\begin{center}\rule{0.5\linewidth}{0.5pt}\end{center}

\subsection{METODOLOGÍA: Cómo Redactar
(3)}\label{metodologuxeda-cuxf3mo-redactar-3}

\textbf{PÁRRAFO 3: TÉCNICAS E INSTRUMENTOS}

{TÉCNICAS}

Se utilizaron tres técnicas principales: (a) entrevistas
semiestructuradas, (b) observación participante en faenas de agua, y (c)
análisis documental de actas comunales.

{INSTRUMENTOS}

Los instrumentos incluyeron guías de entrevista con 15 preguntas
abiertas, fichas de observación estructurada con 8 dimensiones, y matriz
de análisis documental.

{VALIDACIÓN}

Los instrumentos fueron validados por juicio de tres expertos en
antropología andina y recursos hídricos, alcanzando un coeficiente V de
Aiken superior a 0.85.

\begin{center}\rule{0.5\linewidth}{0.5pt}\end{center}

\subsection{METODOLOGÍA: Cómo Redactar
(4)}\label{metodologuxeda-cuxf3mo-redactar-4}

\textbf{PÁRRAFO 4: PROCEDIMIENTO}

{FASE 1}

Coordinación con autoridades comunales y obtención de consentimiento
informado (enero 2023).

{FASE 2}

Trabajo de campo: aplicación de entrevistas, observación de faenas y
rituales, revisión de documentos (febrero-julio 2023).

{FASE 3}

Análisis de datos mediante codificación temática en software Atlas.ti,
identificando categorías emergentes y patrones recurrentes
(agosto-octubre 2023).

\begin{center}\rule{0.5\linewidth}{0.5pt}\end{center}

\subsection{RESULTADOS: Cómo
Redactar}\label{resultados-cuxf3mo-redactar}

\textbf{PRINCIPIOS BÁSICOS}

{ORGANIZACIÓN}

Presenta los hallazgos en el mismo orden que los objetivos específicos.

{OBJETIVIDAD}

Describe lo encontrado sin interpretarlo (la interpretación va en
Discusión).

{EVIDENCIA}

Usa tablas, gráficos, citas textuales de entrevistas para sustentar tus
afirmaciones.

{TIEMPO VERBAL}

Redacta en tiempo pasado: ``Se encontró\ldots{}'', ``Los datos
mostraron\ldots{}'', ``El 65\% manifestó\ldots{}''

\begin{center}\rule{0.5\linewidth}{0.5pt}\end{center}

\subsection{RESULTADOS: Ejemplo de
Redacción}\label{resultados-ejemplo-de-redacciuxf3n}

\textbf{4.1 Conocimientos tradicionales vigentes sobre gestión del agua}

{HALLAZGO PRINCIPAL}

Se identificaron cinco sistemas de conocimiento tradicional actualmente
vigentes en Chuschi: (a) calendario hídrico ritual, (b) sistema de
turnos, (c) señas naturales para predicción, (d) técnicas de
almacenamiento, y (e) normas consuetudinarias de distribución.

{DATOS ESPECÍFICOS}

El 87\% de entrevistados (n=28) reconoció el calendario ritual asociado
al agua, mientras que el 100\% afirmó participar en el sistema de
turnos. La Tabla 3 presenta la distribución de conocimientos por grupos
de edad.

{EVIDENCIA CUALITATIVA}

Un yaku alcalde señaló: ``El agua tiene su tiempo, su mes. Nosotros
sabemos cuándo va a venir, cuándo va a faltar. Los abuelos nos enseñaron
a mirar las nubes, los pájaros'' (Entrevista 07, 45 años).

\begin{center}\rule{0.5\linewidth}{0.5pt}\end{center}

\subsection{DISCUSIÓN: Cómo
Redactar}\label{discusiuxf3n-cuxf3mo-redactar}

\textbf{ESTRUCTURA DEL CAPÍTULO}

{INTERPRETACIÓN}

¿Qué significan tus resultados? ¿Por qué obtuviste estos hallazgos?

{COMPARACIÓN}

¿Cómo se relacionan tus hallazgos con la literatura revisada?
¿Confirman, contradicen o matizan estudios previos?

{LIMITACIONES}

¿Qué factores pudieron influir en tus resultados? ¿Qué aspectos no
pudiste abordar?

{IMPLICACIONES}

¿Qué implicaciones tienen tus hallazgos para la teoría, la práctica o
las políticas?

\begin{center}\rule{0.5\linewidth}{0.5pt}\end{center}

\subsection{DISCUSIÓN: Ejemplo de
Redacción}\label{discusiuxf3n-ejemplo-de-redacciuxf3n}

\textbf{5.1 Interpretación de la persistencia de conocimientos
tradicionales}

{INTERPRETACIÓN}

La alta vigencia de conocimientos tradicionales (87-100\% según
dimensión) sugiere que estos sistemas no son meros ``residuos
culturales'', sino estrategias adaptativas que continúan siendo
funcionalmente efectivas en el contexto local.

{COMPARACIÓN LITERATURA}

Este hallazgo confirma la propuesta de Boelens (2008) sobre la
resiliencia de sistemas consuetudinarios andinos, pero matiza su
argumento al mostrar que la persistencia no es homogénea: algunos
conocimientos (turnos, calendario) son universales, mientras otros
(señas naturales) se concentran en generaciones mayores.

{EXPLICACIÓN}

La diferenciación generacional podría explicarse por el acceso
diferencial a educación formal: los jóvenes con estudios técnicos
tienden a priorizar indicadores meteorológicos formales sobre señas
naturales tradicionales.

\begin{center}\rule{0.5\linewidth}{0.5pt}\end{center}

\subsection{CONCLUSIONES: Cómo Redactar
(Investigación)}\label{conclusiones-cuxf3mo-redactar-investigaciuxf3n}

\textbf{PÁRRAFO 1: RESPUESTA A OBJETIVOS}

{OBJETIVO GENERAL}

La investigación demostró que en Chuschi existe un proceso de
articulación compleja entre conocimientos tradicionales y prácticas
contemporáneas, caracterizado por hibridación selectiva antes que por
reemplazo o conflicto absoluto.

{OBJETIVO ESPECÍFICO 1}

Se identificaron cinco sistemas de conocimiento tradicional vigentes,
con diferentes grados de transmisión intergeneracional.

{OBJETIVO ESPECÍFICO 2}

Las prácticas contemporáneas introducidas incluyen tecnologías de riego
tecnificado, pero su adopción está mediada por criterios de
compatibilidad con la organización comunal.

\begin{center}\rule{0.5\linewidth}{0.5pt}\end{center}

\subsection{CONCLUSIONES: Cómo Redactar (Investigación)
(2)}\label{conclusiones-cuxf3mo-redactar-investigaciuxf3n-2}

\textbf{PÁRRAFO 2: RECOMENDACIONES}

{PARA LA PRÁCTICA}

Se recomienda que los programas de intervención en gestión hídrica
incorporen espacios de diálogo intercultural que permitan integrar ambos
sistemas de conocimiento, evitando imposiciones tecnológicas
unidireccionales.

{PARA LAS POLÍTICAS}

Las políticas públicas de recursos hídricos deberían reconocer
formalmente los sistemas consuetudinarios como formas legítimas de
gobernanza del agua, no solo como ``usos y costumbres'' subordinados a
la legislación estatal.

{PARA LA INVESTIGACIÓN}

Futuras investigaciones podrían profundizar en los mecanismos
específicos de toma de decisiones comunales ante propuestas de
innovación tecnológica, mediante estudios de caso longitudinales.

\begin{center}\rule{0.5\linewidth}{0.5pt}\end{center}

\section{ESTRUCTURA IMRyD}\label{estructura-imryd}

\subsection{El Formato IMRyD}\label{el-formato-imryd}

\subsubsection{¿Qué es IMRyD?}\label{quuxe9-es-imryd}

\textbf{I}ntroducción\\
\textbf{M}étodos\\
\textbf{R}esultados\\
\textbf{y} (and)\\
\textbf{D}iscusión

\subsubsection{Variante Extendida}\label{variante-extendida}

\textbf{I}(AT)\textbf{MRyD}\\
- \textbf{A}ntecedentes\\
- \textbf{T}eorema/Marco Teórico

\subsubsection{Estructura de ``Reloj de
Arena''}\label{estructura-de-reloj-de-arena}

\begin{verbatim}
    General (Introducción)
         ↓
    Específico (Métodos)
         ↓
    Específico (Resultados)
         ↓
    General (Discusión)
\end{verbatim}

\subsubsection{Ventaja}\label{ventaja}

Permite a especialistas \textbf{localizar rápidamente} la información
que necesitan.

\begin{center}\rule{0.5\linewidth}{0.5pt}\end{center}

\subsection{Título Atractivo}\label{tuxedtulo-atractivo}

\subsubsection{Función del Título}\label{funciuxf3n-del-tuxedtulo}

\begin{itemize}
\tightlist
\item
  \textbf{Envoltorio del producto}
\item
  \textbf{Imán} para lectores
\item
  Primera impresión crítica
\item
  Permite decidir pertinencia
\end{itemize}

\subsubsection{Características}\label{caracteruxedsticas}

\begin{itemize}
\tightlist
\item
  Corto (\textless{} 120 caracteres / \textasciitilde14 palabras)\\
\item
  Sencillo y claro\\
\item
  Fiel al contenido\\
\item
  Provocador (opcional)\\
\item
  Sin abreviaturas/siglas
\end{itemize}

\subsubsection{Tipos}\label{tipos}

\begin{enumerate}
\def\labelenumi{\arabic{enumi}.}
\tightlist
\item
  \textbf{Descriptivo:} Qué se estudió

  \begin{itemize}
  \tightlist
  \item
    \emph{``Efecto del fuego sobre la diversidad de gramíneas''}
  \end{itemize}
\item
  \textbf{Informativo:} Resultado principal

  \begin{itemize}
  \tightlist
  \item
    \emph{``El fuego incrementa la diversidad de gramíneas''}
  \end{itemize}
\end{enumerate}

\subsubsection{Evitar}\label{evitar}

\begin{itemize}
\tightlist
\item
  Frases poco informativas:

  \begin{itemize}
  \tightlist
  \item
    ``Estudios sobre\ldots{}''
  \item
    ``Observaciones de\ldots{}''
  \item
    ``Investigaciones acerca de\ldots{}''
  \end{itemize}
\item
  Exceso de detalle técnico
\item
  Limitación geográfica excesiva
\item
  Nombres científicos complejos
\end{itemize}

\subsubsection{Estrategias Creativas}\label{estrategias-creativas}

\textbf{Título en dos partes} (Metáfora : Información)

Ejemplos: - \emph{``De regreso al futuro: Evolución del cerebro
humano''} - \emph{``¿Son mejores los cerebros más grandes?''}

\subsubsection{Regla de Oro}\label{regla-de-oro}

\begin{quote}
Un título atractivo debe \textbf{prometer} básicamente lo que
\textbf{muestra}.
\end{quote}

\begin{center}\rule{0.5\linewidth}{0.5pt}\end{center}

\subsection{Resumen (Abstract)}\label{resumen-abstract}

\subsubsection{El Abstract como
``Asamblea''}\label{el-abstract-como-asamblea}

\begin{quote}
Versión \textbf{en miniatura} del artículo que sintetiza propósito,
métodos, resultados y conclusiones.
\end{quote}

\subsubsection{Estructura (Patrón de Reloj de
Arena)}\label{estructura-patruxf3n-de-reloj-de-arena}

\begin{enumerate}
\def\labelenumi{\arabic{enumi}.}
\tightlist
\item
  \textbf{Introducción mini} → Propósito y contexto
\item
  \textbf{Métodos mini} → Procedimientos generales
\item
  \textbf{Resultados mini} → Hallazgos principales
\item
  \textbf{Discusión mini} → Conclusiones e implicaciones
\end{enumerate}

\subsubsection{Tipos}\label{tipos-1}

\begin{itemize}
\tightlist
\item
  \textbf{Informativo:} Comunica resultados y conclusiones (recomendado)
\item
  \textbf{Descriptivo:} Solo menciona tema y propósito
\end{itemize}

\subsubsection{Requisitos Formales}\label{requisitos-formales}

\begin{itemize}
\tightlist
\item
  \textbf{Extensión:} 150-300 palabras
\item
  \textbf{Formato:} Un solo párrafo
\item
  \textbf{Tiempo verbal:} Pasado
\item
  \textbf{Idiomas:} Original + inglés
\item
  \textbf{Sin:} Citas, referencias a tablas/figuras
\end{itemize}

\subsubsection{Errores Comunes}\label{errores-comunes-1}

\begin{itemize}
\tightlist
\item
  Falta de claridad en mensaje principal\\
\item
  Datos inconsistentes con el texto\\
\item
  Conclusiones no justificadas\\
\item
  Omitir descubrimientos importantes
\end{itemize}

\subsubsection{Importancia}\label{importancia}

Probabilidad de lectura: - \textbf{Solo título + abstract:} 10x -
\textbf{Artículo completo:} 1x

\begin{center}\rule{0.5\linewidth}{0.5pt}\end{center}

\subsection{Introducción:
Justificación}\label{introducciuxf3n-justificaciuxf3n}

\subsubsection{Modelo CARS}\label{modelo-cars}

\textbf{C}reate \textbf{A} \textbf{R}esearch \textbf{S}pace

\paragraph{1. Establecer un Territorio}\label{establecer-un-territorio}

\begin{itemize}
\tightlist
\item
  Demostrar importancia del área
\item
  Introducir investigaciones previas
\item
  Marco conceptual general
\end{itemize}

\paragraph{2. Establecer un Nicho}\label{establecer-un-nicho}

\begin{itemize}
\tightlist
\item
  Señalar vacío de conocimiento
\item
  Plantear interrogante no resuelta
\item
  Extender línea de conocimiento
\end{itemize}

\paragraph{3. Ocupar el Nicho}\label{ocupar-el-nicho}

\begin{itemize}
\tightlist
\item
  Resumir cómo se llenará el vacío
\item
  Objetivos específicos
\item
  Hipótesis (si aplica)
\end{itemize}

\subsubsection{Patrón del Embudo}\label{patruxf3n-del-embudo}

\begin{verbatim}
  ╔═══════════════════╗
  ║  Problema General ║
  ╚═════════╦═════════╝
            ║
       ╔════╩════╗
       ║ Contexto ║
       ╚════╦════╝
            ║
        ╔═══╩═══╗
        ║ Vacío  ║
        ╚═══╦═══╝
            ║
          ╔═╩═╗
          ║ ? ║  ← Pregunta específica
          ╚═══╝
\end{verbatim}

\subsubsection{Extensión Recomendada}\label{extensiuxf3n-recomendada}

\begin{itemize}
\tightlist
\item
  1,000 palabras máximo
\item
  3-4 párrafos bien estructurados
\item
  Solo antecedentes necesarios
\end{itemize}

\begin{center}\rule{0.5\linewidth}{0.5pt}\end{center}

\subsection{Métodos y Materiales}\label{muxe9todos-y-materiales}

\subsubsection{Objetivo Principal}\label{objetivo-principal}

\begin{quote}
Proveer \textbf{detalles suficientes} para que la investigación pueda
ser \textbf{replicada} por otros científicos.
\end{quote}

\subsubsection{Preguntas Clave}\label{preguntas-clave}

\begin{itemize}
\tightlist
\item
  \textbf{¿DÓNDE?} → Lugar del estudio
\item
  \textbf{¿CUÁNDO?} → Período temporal
\item
  \textbf{¿CÓMO?} → Procedimientos exactos
\end{itemize}

\subsubsection{Contenido}\label{contenido}

\begin{itemize}
\tightlist
\item
  Diseño experimental\\
\item
  Materiales utilizados\\
\item
  Técnicas de recolección\\
\item
  Análisis estadísticos\\
\item
  Consideraciones éticas\\
\end{itemize}

\subsubsection{Características de
Redacción}\label{caracteruxedsticas-de-redacciuxf3n}

\begin{itemize}
\tightlist
\item
  \textbf{Tiempo verbal:} Pasado
\item
  \textbf{Voz:} Preferentemente activa
\item
  \textbf{Detalle:} Técnico y preciso
\item
  \textbf{Orden:} Cronológico o lógico
\end{itemize}

\subsubsection{Nivel de Detalle}\label{nivel-de-detalle}

\begin{tcolorbox}[enhanced jigsaw, arc=.35mm, rightrule=.15mm, breakable, titlerule=0mm, title=\textcolor{quarto-callout-tip-color}{\faLightbulb}\hspace{0.5em}{Regla Práctica}, toprule=.15mm, left=2mm, colback=white, colbacktitle=quarto-callout-tip-color!10!white, bottomrule=.15mm, colframe=quarto-callout-tip-color-frame, toptitle=1mm, opacityback=0, opacitybacktitle=0.6, coltitle=black, bottomtitle=1mm, leftrule=.75mm]

Si un colega de tu campo \textbf{NO puede} repetir tu estudio leyendo
esta sección, falta información crítica.

\end{tcolorbox}

\subsubsection{Elementos a Incluir}\label{elementos-a-incluir}

\begin{itemize}
\tightlist
\item
  Nombre científico de especies
\item
  Marcas/modelos de equipos
\item
  Concentraciones exactas
\item
  Tamaños de muestra
\item
  Software utilizado (versión)
\end{itemize}

\begin{center}\rule{0.5\linewidth}{0.5pt}\end{center}

\subsection{Resultados: Presentación
Objetiva}\label{resultados-presentaciuxf3n-objetiva}

\subsubsection{Naturaleza de la
Sección}\label{naturaleza-de-la-secciuxf3n}

\begin{itemize}
\tightlist
\item
  \textbf{Más objetiva} de todo el manuscrito
\item
  Presentación \textbf{sin interpretación}
\item
  No comparar con otros estudios
\item
  No explicar causas
\end{itemize}

\subsubsection{Elementos de
Presentación}\label{elementos-de-presentaciuxf3n}

\begin{enumerate}
\def\labelenumi{\arabic{enumi}.}
\tightlist
\item
  \textbf{Texto narrativo}

  \begin{itemize}
  \tightlist
  \item
    Hallazgos principales
  \item
    Patrones observados
  \end{itemize}
\item
  \textbf{Tablas}

  \begin{itemize}
  \tightlist
  \item
    Datos numéricos precisos
  \item
    Comparaciones múltiples
  \end{itemize}
\item
  \textbf{Figuras/Gráficos}

  \begin{itemize}
  \tightlist
  \item
    Tendencias visuales
  \item
    Patrones biológicos
  \end{itemize}
\end{enumerate}

\subsubsection{Buenas Prácticas}\label{buenas-pruxe1cticas}

\begin{itemize}
\tightlist
\item
  Presentar datos en orden lógico\\
\item
  Destacar hallazgos importantes\\
\item
  Referirse a tablas/figuras\\
\item
  Usar subtítulos si es extenso
\end{itemize}

\subsubsection{Evitar}\label{evitar-1}

\begin{itemize}
\tightlist
\item
  Repetir en texto lo que ya está en tablas\\
\item
  Incluir datos no mencionados en Métodos\\
\item
  Interpretar o especular\\
\item
  Citar literatura (va en Discusión)
\end{itemize}

\subsubsection{Tiempo Verbal}\label{tiempo-verbal}

\begin{itemize}
\tightlist
\item
  \textbf{Pasado:} ``Se observó\ldots{}'', ``Los datos
  mostraron\ldots{}''
\item
  Describe lo que YA ocurrió en TU estudio
\end{itemize}

\begin{center}\rule{0.5\linewidth}{0.5pt}\end{center}

\subsection{Discusión: Interpretación y
Cierre}\label{discusiuxf3n-interpretaciuxf3n-y-cierre}

\subsubsection{La Sección Más
Subjetiva}\label{la-secciuxf3n-muxe1s-subjetiva}

\begin{quote}
Aquí el autor \textbf{interpreta} hallazgos, \textbf{explica} causas y
\textbf{compara} con conocimiento previo.
\end{quote}

\subsubsection{Estructura Recomendada}\label{estructura-recomendada}

\begin{enumerate}
\def\labelenumi{\arabic{enumi}.}
\tightlist
\item
  \textbf{Resumen breve} de hallazgos principales
\item
  \textbf{Interpretación} de resultados
\item
  \textbf{Comparación} con literatura
\item
  \textbf{Explicación} de diferencias/similitudes
\item
  \textbf{Limitaciones} del estudio
\item
  \textbf{Implicaciones} teóricas/prácticas
\item
  \textbf{Futuras investigaciones}
\end{enumerate}

\subsubsection{Relación con
Introducción}\label{relaciuxf3n-con-introducciuxf3n}

La Discusión debe \textbf{cerrar el bucle}: - Introducción: General →
Específico - Discusión: Específico → General

\subsubsection{Voz del Autor}\label{voz-del-autor}

La subjetividad es aceptable y necesaria: - Proponer interpretaciones -
Emitir opiniones fundamentadas - Autocrítica (limitaciones)

\subsubsection{Errores Comunes}\label{errores-comunes-2}

\begin{itemize}
\tightlist
\item
  Repetir Resultados sin interpretar\\
\item
  Introducir datos nuevos\\
\item
  Especular sin fundamento\\
\item
  Ignorar resultados contradictorios
\end{itemize}

\subsubsection{Arte de Convencer}\label{arte-de-convencer}

\begin{quote}
El éxito radica en \textbf{argumentar} de tal manera que el lector
acepte tu interpretación como la más adecuada.
\end{quote}

\subsubsection{Honestidad}\label{honestidad}

Reconocer limitaciones proyecta \textbf{integridad profesional}.

\begin{center}\rule{0.5\linewidth}{0.5pt}\end{center}

\section{PROCESO DE COMPOSICIÓN}\label{proceso-de-composiciuxf3n}

\subsection{Escribir NO es Transvasar
Ideas}\label{escribir-no-es-transvasar-ideas}

\subsubsection{Mito}\label{mito}

\begin{quote}
``Escribir es vaciar el cerebro al papel''
\end{quote}

\subsubsection{Realidad}\label{realidad}

\begin{quote}
Escribir es un \textbf{proceso creativo complejo} que requiere
\textbf{reflexión} sobre el lenguaje
\end{quote}

\subsubsection{Dos Modelos}\label{dos-modelos}

\begin{enumerate}
\def\labelenumi{\arabic{enumi}.}
\tightlist
\item
  \textbf{Contar lo que sé}

  \begin{itemize}
  \tightlist
  \item
    Lineal y simple
  \item
    Vuelco directo de información
  \end{itemize}
\item
  \textbf{Transformar el conocimiento}

  \begin{itemize}
  \tightlist
  \item
    Aprender escribiendo
  \item
    Reorganizar pensamientos
  \item
    Encontrar mejores argumentos
  \end{itemize}
\end{enumerate}

\subsubsection{Proceso No Lineal}\label{proceso-no-lineal}

\begin{verbatim}
Precomposición
     ↓
Redacción ←→ Revisión
     ↓         ↑
Redacción ←→ Revisión
     ↓         ↑
  Resultado
\end{verbatim}

\subsubsection{Beneficio}\label{beneficio}

\begin{tcolorbox}[enhanced jigsaw, arc=.35mm, rightrule=.15mm, breakable, titlerule=0mm, title=\textcolor{quarto-callout-important-color}{\faExclamation}\hspace{0.5em}{Crecimiento Cognitivo}, toprule=.15mm, left=2mm, colback=white, colbacktitle=quarto-callout-important-color!10!white, bottomrule=.15mm, colframe=quarto-callout-important-color-frame, toptitle=1mm, opacityback=0, opacitybacktitle=0.6, coltitle=black, bottomtitle=1mm, leftrule=.75mm]

Al final del proceso, \textbf{sabes MÁS} sobre el tema que cuando
empezaste.

\end{tcolorbox}

\begin{center}\rule{0.5\linewidth}{0.5pt}\end{center}

\subsection{Precomposición: El Paso
Previo}\label{precomposiciuxf3n-el-paso-previo}

\subsubsection{¿Qué es?}\label{quuxe9-es}

\begin{quote}
Etapa de \textbf{planificación} antes de escribir la primera línea.
\end{quote}

\subsubsection{3 Actividades
Simultáneas}\label{actividades-simultuxe1neas}

\paragraph{1. Toma de Conciencia}\label{toma-de-conciencia}

\begin{itemize}
\tightlist
\item
  Reflexión sobre el acto de escribir
\item
  Identificar género y función
\item
  Delimitar tema
\item
  Considerar audiencia
\end{itemize}

\paragraph{2. Descubrimiento}\label{descubrimiento}

\begin{itemize}
\tightlist
\item
  Exploración de memoria
\item
  Asociación libre de ideas
\item
  Generar conexiones nuevas
\item
  Mapas mentales
\end{itemize}

\paragraph{3. Investigación}\label{investigaciuxf3n}

\begin{itemize}
\tightlist
\item
  Acopio de información externa
\item
  Consulta bibliográfica
\item
  Entrevistas/experimentos
\item
  Dominar el tema
\end{itemize}

\subsubsection{Herramienta: Mapa de
Ideas}\label{herramienta-mapa-de-ideas}

\begin{verbatim}
        TEMA CENTRAL
       /     |     \
      /      |      \
  Idea A  Idea B  Idea C
   /  \     |      /  \
 ...  ... ...   ... ...
\end{verbatim}

\textbf{Transforma} pensamiento no lineal → discurso lineal

\subsubsection{Meta}\label{meta}

Pasar de \textbf{escribir para uno mismo} a \textbf{escribir para el
lector}.

\begin{center}\rule{0.5\linewidth}{0.5pt}\end{center}

\subsection{Identificación del
Público}\label{identificaciuxf3n-del-puxfablico}

\subsubsection{``Cálculo del Lector''}\label{cuxe1lculo-del-lector}

\begin{quote}
Prever \textbf{reacciones}, \textbf{conocimientos} y
\textbf{necesidades} del destinatario.
\end{quote}

\subsubsection{Preguntas}\label{preguntas}

\begin{itemize}
\tightlist
\item
  ¿Quiénes son mis lectores reales?
\item
  ¿Qué nivel cultural tienen?
\item
  ¿Cuánto saben del tema?
\item
  ¿Cuánto tiempo disponen para leer?
\item
  ¿Por qué les interesaría?
\end{itemize}

\subsubsection{Dos Figuras Discursivas}\label{dos-figuras-discursivas}

\begin{enumerate}
\def\labelenumi{\arabic{enumi}.}
\tightlist
\item
  \textbf{Autor textual}

  \begin{itemize}
  \tightlist
  \item
    Imagen que proyecto
  \item
    Rigor, inteligencia, credibilidad
  \end{itemize}
\item
  \textbf{Lector textual}

  \begin{itemize}
  \tightlist
  \item
    Destinatario ideal construido
  \item
    Basado en lectores reales
  \end{itemize}
\end{enumerate}

\subsubsection{Registro y Tenor}\label{registro-y-tenor}

El público determina: - Nivel de formalidad - Uso de tecnicismos -
Cantidad de explicaciones - Tono (tú vs.~usted)

\subsubsection{Lo Explícito
vs.~Implícito}\label{lo-expluxedcito-vs.-impluxedcito}

Decisión crítica: - ¿Cuánta información dar? - ¿Qué dejar a inferencias?

\subsubsection{Errores}\label{errores}

\begin{itemize}
\tightlist
\item
  \textbf{Exceso:} Tratar al lector de ignorante\\
\item
  \textbf{Defecto:} Omitir datos necesarios
\end{itemize}

\subsubsection{Objetivo}\label{objetivo}

Convertir al lector real en \textbf{interlocutor necesario} para
comunicación eficaz.

\begin{center}\rule{0.5\linewidth}{0.5pt}\end{center}

\subsection{Redacción: De lo No Lineal a lo
Lineal}\label{redacciuxf3n-de-lo-no-lineal-a-lo-lineal}

\subsubsection{Resolución de
Problemas}\label{resoluciuxf3n-de-problemas}

Escribir implica decidir simultáneamente: - \textbf{Qué} decir -
\textbf{Cómo} decirlo - \textbf{Qué voz} adoptar - \textbf{Qué léxico}
usar

\subsubsection{Problemas Circulares}\label{problemas-circulares}

Una solución estilística puede generar un nuevo problema conceptual.

\subsubsection{Fórmula 18/50}\label{fuxf3rmula-1850}

\begin{itemize}
\tightlist
\item
  \textbf{18 palabras} por oración\\
\item
  \textbf{50 palabras} por párrafo (\textasciitilde3 oraciones)
\end{itemize}

\subsubsection{Voz Activa \textgreater{}
Pasiva}\label{voz-activa-pasiva}

\begin{itemize}
\tightlist
\item
  \textbf{Activa:} ``El investigador analizó los datos''\\
\item
  \textbf{Pasiva:} ``Los datos fueron analizados''
\end{itemize}

\subsubsection{Principios de
Redacción}\label{principios-de-redacciuxf3n}

\begin{enumerate}
\def\labelenumi{\arabic{enumi}.}
\tightlist
\item
  \textbf{Claridad y sencillez}

  \begin{itemize}
  \tightlist
  \item
    Evitar ``lenguaje infiltrado''
  \item
    Prosa llana y directa
  \end{itemize}
\item
  \textbf{Precisión léxica}

  \begin{itemize}
  \tightlist
  \item
    Palabras cortas y conocidas
  \item
    Evitar sinónimos inexactos
  \end{itemize}
\item
  \textbf{Orden lógico}

  \begin{itemize}
  \tightlist
  \item
    Sujeto + Verbo + Objeto
  \end{itemize}
\end{enumerate}

\subsubsection{En Textos Científicos
(IMRyD)}\label{en-textos-cientuxedficos-imryd}

\begin{itemize}
\tightlist
\item
  \textbf{Impersonal} y objetivo
\item
  Hallazgos publicados: \textbf{presente}
\item
  Resultados propios: \textbf{pasado}
\end{itemize}

\subsubsection{Ejemplo}\label{ejemplo}

\begin{itemize}
\tightlist
\item
  ``Se observó un aumento del 23\%''\\
\item
  ``Hubo un incremento significativo''
\end{itemize}

\begin{center}\rule{0.5\linewidth}{0.5pt}\end{center}

\subsection{Borradores Sucesivos}\label{borradores-sucesivos}

\subsubsection{Principio Fundamental}\label{principio-fundamental}

\begin{quote}
\textbf{Escribir y reescribir son la misma cosa.}
\end{quote}

\subsubsection{¿Por qué son
necesarios?}\label{por-quuxe9-son-necesarios}

\begin{itemize}
\tightlist
\item
  El lenguaje es \textbf{ambiguo} por naturaleza
\item
  Primera versión rara vez es óptima
\item
  Ideas maduran con revisión
\item
  Claridad requiere ``afilar'' el texto
\end{itemize}

\subsubsection{Proceso de Revisión}\label{proceso-de-revisiuxf3n}

\begin{enumerate}
\def\labelenumi{\arabic{enumi}.}
\tightlist
\item
  \textbf{Escribir} primer borrador
\item
  \textbf{Dejar reposar} el texto
\item
  \textbf{Releer} con ojos frescos
\item
  \textbf{Tachar} lo superfluo
\item
  \textbf{Reorganizar} párrafos
\item
  \textbf{Someter} a lectura ajena
\item
  \textbf{Repetir} hasta satisfacción
\end{enumerate}

\subsubsection{Meta}\label{meta-1}

Crear la \textbf{ilusión} de que el texto se escribió de una sola vez,
con fluidez.

\subsubsection{Técnica del Tachado}\label{tuxe9cnica-del-tachado}

Preguntas durante revisión: - ¿Puedo decir esto con menos palabras? -
¿Esta oración añade valor? - ¿Hay redundancia? - ¿El orden es lógico?

\subsubsection{Límite del
Perfeccionamiento}\label{luxedmite-del-perfeccionamiento}

Llega un punto donde el texto se vuelve ``intratable'' → aceptar como
definitivo.

\subsubsection{Lectura Ajena}\label{lectura-ajena}

\begin{tcolorbox}[enhanced jigsaw, arc=.35mm, rightrule=.15mm, breakable, titlerule=0mm, title=\textcolor{quarto-callout-tip-color}{\faLightbulb}\hspace{0.5em}{Consejo de Oro}, toprule=.15mm, left=2mm, colback=white, colbacktitle=quarto-callout-tip-color!10!white, bottomrule=.15mm, colframe=quarto-callout-tip-color-frame, toptitle=1mm, opacityback=0, opacitybacktitle=0.6, coltitle=black, bottomtitle=1mm, leftrule=.75mm]

Incluso borradores ``finales'' se benefician de otra perspectiva para
identificar pasajes oscuros.

\end{tcolorbox}

\begin{center}\rule{0.5\linewidth}{0.5pt}\end{center}

\subsection{La oración}\label{la-oraciuxf3n}

\subsubsection{Longitud de la Oración: Índice de
Nebulosidad}\label{longitud-de-la-oraciuxf3n-uxedndice-de-nebulosidad}

\textbf{Fórmula de Fernando Ávila (2003)}

\[\text{Índice de Nebulosidad} = (A + B) \times 0.4\]

Donde:

\begin{itemize}
\tightlist
\item
  \textbf{A} = Número de palabras / Número de oraciones
\item
  \textbf{B} = (Palabras largas × 100) / Número total de palabras
\end{itemize}

\textbf{Índice ideal}: 10 (nivel ``Selecciones'')

\begin{center}\rule{0.5\linewidth}{0.5pt}\end{center}

\subsubsection{Ejemplo 1: Oración Larga
Problemática}\label{ejemplo-1-oraciuxf3n-larga-problemuxe1tica}

\textbf{VERSIÓN ORIGINAL} (134 palabras, 1 oración)

De acuerdo con especialistas de nuestro departamento de calidad en
metodología internacional para la presentación de trabajos escritos
debidamente asesorados por una comisión de la Universidad Autónoma de la
Nación, el texto de la novela con la que usted puede participar en el
Decimonoveno Concurso Nacional de Narrativa, patrocinado por el
Ministerio de Cultura, la Fundación para el Desarrollo de la
Inteligencia Artística Continental y el grupo de editoriales educativas
legalmente afincadas en el territorio de la patria, debe estar escrito
originalmente en castellano, en letra cuerpo 12, sobre hojas de papel
bond base 20 tamaño carta, a doble espacio, con sangría de cuatro
espacios.

\textbf{Problemas identificados}: Una sola oración, múltiples ideas,
difícil seguimiento

\begin{center}\rule{0.5\linewidth}{0.5pt}\end{center}

\subsubsection{Ejemplo 1: Versión
Mejorada}\label{ejemplo-1-versiuxf3n-mejorada}

\textbf{VERSIÓN MEJORADA} (3 oraciones)

Los especialistas de nuestro departamento han establecido los requisitos
formales para el Decimonoveno Concurso Nacional de Narrativa.
\textbf{{[}Oración 1: Contexto{]}}

El texto debe estar escrito originalmente en castellano, en letra cuerpo
12, sobre hojas de papel bond base 20 tamaño carta. \textbf{{[}Oración
2: Requisitos técnicos{]}}

Además, debe presentarse a doble espacio con sangría de cuatro espacios.
\textbf{{[}Oración 3: Formato específico{]}}

\textbf{Mejoras logradas}: Claridad, separación de ideas, lectura fluida

\begin{center}\rule{0.5\linewidth}{0.5pt}\end{center}

\subsubsection{Estructura Oracional: Quién + Verbo + Qué +
Complemento}\label{estructura-oracional-quiuxe9n-verbo-quuxe9-complemento}

\textbf{QUIÉN} Sujeto Agente

\textbf{VERBO} Acción Núcleo

\textbf{QUÉ} Objeto directo Contenido

\textbf{COMPLEMENTO} A quién, dónde,cuándo, cómo

\textbf{Ejemplo aplicado:}

{\textbf{Diego Pérez}} {\textbf{ofreció}} {\textbf{lo mejor de su
repertorio}} {\textbf{a las jóvenes representantes}}

\begin{center}\rule{0.5\linewidth}{0.5pt}\end{center}

\subsubsection{Ejemplo 2: Reordenamiento de
Estructura}\label{ejemplo-2-reordenamiento-de-estructura}

\textbf{Versión original} (orden confuso):

Entre el primero y el 31 de octubre del presente año, diseñado con gran
esmero por una comisión conjunta, en el próximo seminario de
actualización profesional pueden participar los egresados de la
Universidad Autónoma.

\textbf{Versión ordenada} (Quién + Verbo + Qué + Complementos):

{\textbf{Los egresados de la Universidad Autónoma}} {\textbf{podrán
participar}} {\textbf{en el próximo seminario de actualización
profesional}} {\textbf{entre el 1 y 31 de octubre de 2022}}.

Una comisión conjunta de egresados, profesores y alumnos diseñó el
seminario con gran esmero.

\begin{center}\rule{0.5\linewidth}{0.5pt}\end{center}

\subsubsection{Voz Activa vs.~Voz
Pasiva}\label{voz-activa-vs.-voz-pasiva}

\textbf{VOZ PASIVA} (Evitar)

El experimento fue realizado por Pedro.

Una nueva propuesta es entregada por Margarita a la Junta Directiva.

Un libro de política es leído por Juan.

\textbf{VOZ ACTIVA} (Preferir)

Pedro realizó el experimento.

Margarita entrega una nueva propuesta a la Junta Directiva.

Juan lee un libro de política.

\textbf{Razón}: La voz activa es más directa, clara y concisa.

\begin{center}\rule{0.5\linewidth}{0.5pt}\end{center}

\subsubsection{Persona Gramatical:
Consistencia}\label{persona-gramatical-consistencia}

\textbf{Problema}: Cambio innecesario de persona

Tenemos el gusto de ofrecerle nuestros servicios. \emph{{[}Primera
persona plural{]}}

Se darán descuentos a los primeros compradores. \emph{{[}Impersonal{]}}

Le estaré informando de las promociones. \emph{{[}Primera persona
singular{]}}

\textbf{Solución}: Mantener la misma persona

Tenemos el gusto de ofrecerle nuestros servicios. \emph{{[}Primera
persona plural{]}}

Daremos descuentos a los primeros cien compradores. \emph{{[}Primera
persona plural{]}}

Le informaremos oportunamente de las nuevas promociones.
\emph{{[}Primera persona plural{]}}

\begin{center}\rule{0.5\linewidth}{0.5pt}\end{center}

\subsubsection{Tiempo Verbal en Textos
Científicos}\label{tiempo-verbal-en-textos-cientuxedficos}

\textbf{PRESENTE} Para conocimiento establecido

El fármaco X \textbf{incide} en las tasas de fertilidad de la especie Y
(Smith et al., 2019).

La teoría de la relatividad \textbf{establece} que\ldots{}

\textbf{PASADO} Para tus hallazgos específicos

\textbf{Examinamos} el efecto del fármaco X en diez individuos.

\textbf{Encontramos} que el 40\% de los participantes\ldots{}

\begin{center}\rule{0.5\linewidth}{0.5pt}\end{center}

\subsection{El Párrafo: Unidad de
Pensamiento}\label{el-puxe1rrafo-unidad-de-pensamiento}

\subsubsection{Definición}\label{definiciuxf3n-2}

\begin{quote}
Bloque constituyente que expresa y argumenta \textbf{una única idea
central}.
\end{quote}

\subsubsection{Estructura del Párrafo
Ideal}\label{estructura-del-puxe1rrafo-ideal}

\begin{enumerate}
\def\labelenumi{\arabic{enumi}.}
\tightlist
\item
  \textbf{Oración temática}

  \begin{itemize}
  \tightlist
  \item
    Presenta el tema
  \item
    Idea u opinión (NO hecho)
  \end{itemize}
\item
  \textbf{Oraciones de apoyo} (2-5)

  \begin{itemize}
  \tightlist
  \item
    Justifican la oración temática
  \item
    Datos, ejemplos, información
  \end{itemize}
\item
  \textbf{Oración concluyente}

  \begin{itemize}
  \tightlist
  \item
    Resumen O conexión al siguiente párrafo
  \end{itemize}
\end{enumerate}

\subsubsection{Extensión Recomendada}\label{extensiuxf3n-recomendada-1}

\begin{itemize}
\tightlist
\item
  \textbf{7-14 líneas} (ideal)\\
\item
  \textbf{3-6 líneas} (ocasional, corto)\\
\item
  \textbf{15-20 líneas} (ocasional, largo)\\
\end{itemize}

\subsubsection{Errores Comunes}\label{errores-comunes-3}

\begin{enumerate}
\def\labelenumi{\arabic{enumi}.}
\tightlist
\item
  Párrafo de \textbf{una sola oración}
\item
  Párrafo que ocupa \textbf{página entera}
\item
  \textbf{Múltiples ideas} en un párrafo
\end{enumerate}

\subsubsection{Soluciones}\label{soluciones}

\textbf{Si no transmite idea:} ELIMINAR\\
\textbf{Si transmite muchas:} FRAGMENTAR

\subsubsection{Principios de Coherencia}\label{principios-de-coherencia}

\begin{itemize}
\tightlist
\item
  \textbf{Unidad:} Todas las oraciones relacionadas con idea principal\\
\item
  \textbf{Orden lógico:} Cronológico, importancia o alcance\\
\item
  \textbf{Independencia relativa:} Puede leerse solo
\end{itemize}

\subsubsection{Variedad}\label{variedad}

Alternar párrafos de longitud media con cortos/largos evita monotonía.

\begin{center}\rule{0.5\linewidth}{0.5pt}\end{center}

\subsubsection{Estructura 1: Afirmación → Explicación →
Ejemplo}\label{estructura-1-afirmaciuxf3n-explicaciuxf3n-ejemplo}

\textbf{Esquema visual:}

\begin{verbatim}
[AFIRMACIÓN GENERAL] ────┐
                          │
[EXPLICACIÓN CONTEXTUAL] ─┼─ DESARROLLO
                          │
[EJEMPLO ESPECÍFICO] ─────┘
\end{verbatim}

\textbf{Ejemplo real} (Xie et al., 2011):

{\textbf{El comportamiento humano está profundamente afectado por la
influencia de individuos y de las redes sociales que los vinculan.}}
Mucho antes de la proliferación de redes sociales on-line, los
estudiosos ya habían reconocido las redes sociales off-line como un
factor importante que determina cómo las sociedades se mueven hacia
consensos. {En la sociología, trabajos sobre la difusión de innovaciones
han enfatizado cómo los individuos adoptan nuevos comportamientos por
influencia de sus vecinos.}

\begin{center}\rule{0.5\linewidth}{0.5pt}\end{center}

\subsubsection{Estructura 2: Cadena Descriptiva → Afirmación
Conclusiva}\label{estructura-2-cadena-descriptiva-afirmaciuxf3n-conclusiva}

\textbf{Esquema visual:}

\begin{verbatim}
DESCRIPCIÓN 1 (contexto físico) ───┐
                                    │
DESCRIPCIÓN 2 (estado inicial) ────┼─ OBSERVACIONES
                                    │
DESCRIPCIÓN 3 (comportamiento) ────┤
                                    ↓
                            AFIRMACIÓN FINAL
                          (significado/función)
\end{verbatim}

\textbf{Ejemplo real} (Dawkins, 1976):

Las gaviotas de cabeza negra hacen sus nidos en grandes colonias, a
pocos metros de distancia. Cuando los polluelos salen del cascarón son
pequeños e indefensos y fáciles de tragar. Es bastante común que una
gaviota espere hasta que un vecino le dé la espalda y entonces se
abalance sobre uno de los polluelos del vecino y se lo trague entero.
{\textbf{Así obtiene una buena comida nutritiva, sin tener que tomarse
la molestia de pescar.}}

\begin{center}\rule{0.5\linewidth}{0.5pt}\end{center}

\subsubsection{Estructura 3: Afirmación Incompleta →
Aclaración}\label{estructura-3-afirmaciuxf3n-incompleta-aclaraciuxf3n}

\textbf{Esquema visual:}

\begin{verbatim}
PARADOJA/PROBLEMA ──→ INTERROGANTE ──→ ACLARACIÓN ──→ RESOLUCIÓN
\end{verbatim}

\textbf{Ejemplo real} (Hernandez, 2011):

{\textbf{Los casos de pérdida y recuperación de lenguas plantean una
intrigante paradoja.}} Si en el cerebro se almacenan dos lenguas, ¿cómo
es posible que una persona pueda perder una lengua y recuperarla? Mi
propia experiencia y los casos de bilingüismo que leía sugerían un
modelo muy diferente. {\textbf{Parecía que las lenguas podían ser
inaccesibles aunque no se perdieran del todo.}}

\begin{center}\rule{0.5\linewidth}{0.5pt}\end{center}

\subsubsection{Estructura 4: Afirmación → Justificación con
Datos}\label{estructura-4-afirmaciuxf3n-justificaciuxf3n-con-datos}

\textbf{Esquema visual:}

\begin{verbatim}
TESIS PRINCIPAL
    ↓
DATO 1 (cuantitativo) ─┐
                        │
DATO 2 (comparativo) ──┼─ EVIDENCIA EMPÍRICA
                        │
DATO 3 (estadístico) ──┘
\end{verbatim}

\textbf{Ejemplo real} (Bertrand et al., 2010):

{\textbf{La presencia de hijos es el principal factor que contribuye a
la menor experiencia laboral de las mujeres con MBA.}} Las mujeres con
hijos tienen un déficit de unos 8 meses de experiencia en comparación
con el hombre medio, mientras que las mujeres sin hijos tienen un
déficit de 1,5 meses. Las mujeres con hijos suelen trabajar un 24\%
menos de horas semanales; las mujeres sin hijos sólo trabajan un 3,3\%
menos.

\begin{center}\rule{0.5\linewidth}{0.5pt}\end{center}

\subsubsection{Estructura 5: Juicio Negativo → Atenuación → Ejemplo
Positivo}\label{estructura-5-juicio-negativo-atenuaciuxf3n-ejemplo-positivo}

\textbf{Esquema visual:}

\begin{verbatim}
CRÍTICA INICIAL (fuerte) ──→ ATENUACIÓN (pero/sin embargo) ──→ CASO EXITOSO
\end{verbatim}

\textbf{Ejemplo real} (Bunge, 2012):

{\textbf{Las hipótesis sobre equilibrio o armonía son tan imprecisas que
no se las puede poner a prueba, de modo que no son científicas.}} Pero
algunas de ellas son pasibles de refinamiento y pueden tornarse
accesibles a la contrastación experimental. {\textbf{En efecto, la
versión moderna de la idea de equilibrio corporal es el concepto
fisiológico de homeostasis, propuesta por Claude Bernard y refinada por
Walter Cannon.}}

\begin{center}\rule{0.5\linewidth}{0.5pt}\end{center}

\subsubsection{Estructura 6: Descripción General → Detalles
Secuenciales}\label{estructura-6-descripciuxf3n-general-detalles-secuenciales}

\textbf{Esquema visual:}

\begin{verbatim}
PRESENTACIÓN GENERAL (contexto + n + muestra)
    ↓
    ├─ HALLAZGO 1 (principal)
    ├─ HALLAZGO 2 (metodología específica)
    ├─ HALLAZGO 3 (implicaciones)
    └─ HALLAZGO 4 (contexto ampliado)
\end{verbatim}

\textbf{Ejemplo real} (Stoet et al., 2018):

{\textbf{Presentamos un análisis del rendimiento académico de casi
475.000 adolescentes de 67 países.}} Descubrimos que las mujeres y los
hombres tienen capacidades similares en conocimientos científicos. Al
mismo tiempo, es mucho más probable que las ciencias sean un punto
fuerte para los chicos que para las chicas. La relación entre estas
diferencias y las tasas de graduación en STEM es mayor en países con
mayor igualdad de género.

\begin{center}\rule{0.5\linewidth}{0.5pt}\end{center}

\subsubsection{Estructura 7: Definición → Explicación → Ejemplo
Contrastivo}\label{estructura-7-definiciuxf3n-explicaciuxf3n-ejemplo-contrastivo}

\textbf{Esquema visual:}

\begin{verbatim}
DEFINICIÓN FORMAL
    ↓
CARACTERÍSTICAS DISTINTIVAS
    ↓
EJEMPLO ILUSTRATIVO ←→ CONTRASTE CON OTROS CASOS
\end{verbatim}

\textbf{Ejemplo real} (Mosterín, 1970):

{\textbf{Los nombres son sucesiones de signos gráficos que designan
algo: un número, una ciudad, una persona.}} Se diferencian por ser más
cortos, más unívocos, más independientes del contexto. Si digo ``yo he
comido'', ``yo'' actúa como designador que se refiere a mí, pero su
referencia varía con el contexto. {\textbf{``Jesús Mosterín'' permanece
invariable, utilice la sentencia quien la utilice.}}

\begin{center}\rule{0.5\linewidth}{0.5pt}\end{center}

\subsubsection{Estructura 8: Interrogante → Respuesta → Justificación
Implícita}\label{estructura-8-interrogante-respuesta-justificaciuxf3n-impluxedcita}

\textbf{Esquema visual:}

\begin{verbatim}
ESCENARIO HIPOTÉTICO (gancho emocional)
    ↓
PREGUNTA FUNDAMENTAL
    ↓
RESPUESTA CONCISA + JUSTIFICACIÓN INTEGRADA
\end{verbatim}

\textbf{Ejemplo real} (Feynman, 1963):

Si en algún cataclismo se destruyera todo el conocimiento científico y
sólo se transmitiera una frase, ¿qué afirmación contendría la mayor
cantidad de información? {\textbf{Creo que es la hipótesis de que todas
las cosas están hechas de átomos: pequeñas partículas en perpetuo
movimiento que se atraen a poca distancia pero se repelen cuando se
aprietan.}}

\begin{center}\rule{0.5\linewidth}{0.5pt}\end{center}

\subsubsection{Estructura 9: Interrogante → Aclaración Metodológica →
Conexión}\label{estructura-9-interrogante-aclaraciuxf3n-metodoluxf3gica-conexiuxf3n}

\textbf{Esquema visual:}

\begin{verbatim}
PREGUNTAS INICIALES (interés)
    ↓
ACLARACIÓN NECESARIA (perspectiva)
    ↓
CONEXIÓN EXPLÍCITA (siguiente párrafo)
    ↓
NUEVA PREGUNTA (momentum)
\end{verbatim}

\textbf{Ejemplo real} (Crystal, 2014):

¿Cuántas lenguas están a punto de morir? ¿Cuántas están en peligro?
{\textbf{Antes de estimar la magnitud del problema, debemos desarrollar
un sentido de perspectiva.}} Las cifras sobre lenguas que mueren sólo
tienen sentido si las relacionamos con el número total de lenguas vivas.
{\textbf{Entonces, ¿cuántas lenguas hay?}}

\begin{center}\rule{0.5\linewidth}{0.5pt}\end{center}

\subsubsection{Características del Párrafo
Efectivo}\label{caracteruxedsticas-del-puxe1rrafo-efectivo}

\textbf{LO QUE SÍ DEBE TENER}

\begin{itemize}
\tightlist
\item
  Coherencia conceptual
\item
  Una idea principal clara
\item
  Extensión moderada (50-60 palabras promedio)
\item
  Conexión lógica interna
\item
  Independencia relativa del contexto
\end{itemize}

\textbf{LO QUE DEBE EVITAR}

\begin{itemize}
\tightlist
\item
  Párrafos de una sola oración
\item
  Párrafos que ocupan página entera
\item
  Múltiples ideas sin conexión
\item
  Dependencia total de otros párrafos
\item
  Más de cuatro oraciones largas
\end{itemize}

\begin{center}\rule{0.5\linewidth}{0.5pt}\end{center}

\subsection{Texto expositivo}\label{texto-expositivo}

\subsubsection{Tipos de Estructura Textual:
Comparativa}\label{tipos-de-estructura-textual-comparativa}

Tipo

Característica

Ejemplo

Temática

Organización por conceptos o categorías

Peñaherrera (2004): Geografía del Perú

Histórica

Organización cronológica pura

Del Busto (2004): Conquista y virreinato

Histórico-temática

Períodos + temas dentro de cada período

Harari (2011): Sapiens

Histórico-crítica

Cronología + análisis crítico

Roll (1938): Historia del pensamiento económico

Mixta

Combina varios principios organizadores

Medrano (2021): Quipus

\begin{center}\rule{0.5\linewidth}{0.5pt}\end{center}

\subsubsection{Estructura Temática: Análisis
Detallado}\label{estructura-temuxe1tica-anuxe1lisis-detallado}

\textbf{Caso:} Peñaherrera (2004) - Geografía del Perú

\textbf{Principio organizador:} De lo general a lo específico, por
categorías geográficas

\begin{verbatim}
1. CONTEXTO AMPLIO
   └─ El espacio nacional en América y el mundo

2. HERRAMIENTAS
   └─ Cartografía peruana

3. PROCESOS FORMATIVOS
   └─ La formación del relieve andino

4. RESULTADOS FÍSICOS
   ├─ Relieve del territorio peruano
   └─ Las regiones naturales del Perú

5. ELEMENTOS ESPECÍFICOS
   ├─ Climas del Perú
   ├─ El mar de Grau
   ├─ Cuencas del Pacífico y del Titicaca
   └─ Cuenca del Amazonas

6. SÍNTESIS HUMANA
   └─ Síntesis de geografía humana
\end{verbatim}

\begin{center}\rule{0.5\linewidth}{0.5pt}\end{center}

\subsubsection{Estructura Histórica: Análisis
Detallado}\label{estructura-histuxf3rica-anuxe1lisis-detallado}

\textbf{Caso:} Del Busto (2004) - Conquista y virreinato

\textbf{Principio organizador:} Cronológico estricto + eventos clave

\textbf{FASE 1: CONQUISTA}

\begin{itemize}
\tightlist
\item
  Los conquistadores españoles
\item
  Los viajes descubridores
\item
  La caída del Tahuantinsuyo
\item
  La fundación de ciudades
\item
  La resistencia incaica
\end{itemize}

\textbf{FASE 2: CONSOLIDACIÓN}

\begin{itemize}
\tightlist
\item
  Las guerras civiles
\item
  Los grandes descubrimientos
\item
  Los incas de Vilcabamba
\end{itemize}

\textbf{FASE 3: VIRREINATO}

\begin{itemize}
\tightlist
\item
  Administración central y virreinal
\item
  Sociedad virreinal
\item
  Catolicismo y religiosidad
\item
  Educación y vida intelectual
\item
  Comercio, industria y artesanía
\item
  Las rebeliones indígenas
\end{itemize}

\begin{center}\rule{0.5\linewidth}{0.5pt}\end{center}

\subsubsection{Estructura Histórico-Temática: Análisis
Detallado}\label{estructura-histuxf3rico-temuxe1tica-anuxe1lisis-detallado}

\textbf{Caso:} Harari (2011) - Sapiens

\textbf{Principio organizador:} Grandes revoluciones + temas dentro de
cada una

Revolución

Temas

Lógica

Cognitiva

Animal insignificante → Árbol conocimiento → Vida Adán y Eva

Biológico → Mental → Social

Agrícola

Mayor fraude → Pirámides → Memoria → Justicia

Cambio → Consecuencias → Problemas

Unificación

Flecha historia → Dinero → Imperios → Religión

Proceso → Herramientas integradoras

Científica

Ignorancia → Imperio → Capitalismo → Revolución permanente

Método → Alianzas → Transformación

\begin{center}\rule{0.5\linewidth}{0.5pt}\end{center}

\subsubsection{Estructura Histórico-Crítica: Análisis
Detallado}\label{estructura-histuxf3rico-cruxedtica-anuxe1lisis-detallado}

\textbf{Caso:} Roll (1938) - Historia del pensamiento económico

\textbf{Principio organizador:} Desarrollo histórico + análisis crítico
de cada escuela

\textbf{FASE 1: Comienzos} → Antiguo Testamento, Grecia, Roma,
Cristianismo, Edad Media

\textbf{FASE 2: Capitalismo comercial} → Declive escolástica,
mercantilismo, lingotismo

\textbf{FASE 3: Fundadores} → Filósofos políticos, Petty, Locke, Hume,
Cantillon, fisiócratas

\textbf{FASE 4: Sistema clásico} → {\textbf{Adam Smith (análisis
extenso)}}, Ricardo, Malthus

\textbf{FASE 5: Reacción} → Crítica a Malthus, románticos alemanes,
crítica socialista

\textbf{FASE 6: Marx} → {\textbf{Vida, método, teoría valor, plusvalía,
competencia, desarrollo, valoración}}

\textbf{FASE 7: Transición} → Escuela histórica, Jones, ruptura teoría
valor, Senior, Mill

\textbf{FASE 8-12: Economía moderna} → Utilidad marginal, contribución
americana, entreguerras, macroeconomía, edad de la duda

\textbf{CONCLUSIÓN:} Los fines de la economía

\begin{center}\rule{0.5\linewidth}{0.5pt}\end{center}

\subsubsection{Estructura Mixta: Análisis
Detallado}\label{estructura-mixta-anuxe1lisis-detallado}

\textbf{Caso:} Medrano (2021) - Quipus

\textbf{Principio organizador:} Histórico (Parte 1) + Analítico
contemporáneo (Parte 2)

\textbf{PARTE 1: HISTÓRICA}

Una historia en nudos

\begin{itemize}
\tightlist
\item
  Quipus tempranos
\item
  Quipus incaicos
\item
  Encuentros coloniales tempranos
\item
  Quipus coloniales y republicanos
\end{itemize}

{Narrativa cronológica tradicional}

\textbf{PARTE 2: CONTEMPORÁNEA}

Desciframientos presentes y futuros

\begin{itemize}
\tightlist
\item
  Hacia el quipu digital
\item
  Los quipus digitales como ``macrodatos''
\end{itemize}

{Análisis metodológico actual}

\textbf{Ventaja de esta estructura:} Combina narrativa accesible
histórica con análisis técnico especializado, permitiendo a diferentes
audiencias encontrar valor en distintas secciones.

\begin{center}\rule{0.5\linewidth}{0.5pt}\end{center}

\subsubsection{Texto Convincente: Elementos
Esenciales}\label{texto-convincente-elementos-esenciales}

\textbf{Según Reyes (p.~259), un texto académico convincente debe:}

\textbf{ARGUMENTACIÓN}

Propone argumentos válidos y suficientes

\textbf{EVIDENCIA}

Provee datos adecuados para el tema y lector

\textbf{INFERENCIA}

Permite que el lector llegue a conclusiones

\textbf{CONEXIÓN}

Conecta datos con la tesis defendida

\textbf{CONTRASTE}

Compara con tesis contrarias

\textbf{APORTE}

Amplía conocimientos del lector

\begin{center}\rule{0.5\linewidth}{0.5pt}\end{center}

\subsubsection{Lista de Verificación de Calidad
Textual}\label{lista-de-verificaciuxf3n-de-calidad-textual}

\textbf{Según Reyes (p.~272):}

\textbf{CLARIDAD Y COHERENCIA}

¿Ha quedado clara la intención del escrito?

¿Se entiende todo bien?

¿Están las conclusiones justificadas?

¿Hay contradicciones?

\textbf{ESTRUCTURA Y FORMA}

¿Es adecuada la estructura?

¿Son acertados registro y vocabulario?

¿Son buenas las frases inicial y final?

¿Está bien contextualizado el problema?

\textbf{REVISIÓN EXTERNA}

Siempre pedir opinión de otras personas: ¿se entiende todo bien?

\begin{center}\rule{0.5\linewidth}{0.5pt}\end{center}

\subsubsection{Aplicación: De Teoría a
Práctica}\label{aplicaciuxf3n-de-teoruxeda-a-pruxe1ctica}

\textbf{Secuencia de redacción recomendada:}

\begin{verbatim}
1. PALABRA (precisa, conocida, apropiada)
        ↓
2. FRASE (corta, correcta, clara)
        ↓
3. ORACIÓN (18 palabras promedio, estructura lógica)
        ↓
4. PÁRRAFO (50 palabras promedio, una idea central)
        ↓
5. TEXTO (estructura clara, argumentación sólida)
\end{verbatim}

\textbf{Recuerde:} Cada nivel debe cumplir criterios específicos de
calidad antes de pasar al siguiente. La revisión debe hacerse en orden
inverso: primero la estructura general del texto, luego los párrafos,
luego las oraciones, y finalmente las palabras individuales.

\begin{center}\rule{0.5\linewidth}{0.5pt}\end{center}

\section{MÓDULO 4: HERRAMIENTAS DE
ESTILO}\label{muxf3dulo-4-herramientas-de-estilo}

\subsection{Vocabulario Preciso: Propiedad
Léxica}\label{vocabulario-preciso-propiedad-luxe9xica}

\subsubsection{Definición}\label{definiciuxf3n-3}

\begin{quote}
Uso de \textbf{palabras exactas} que comunican lo que se desea decir,
sin ambigüedades.
\end{quote}

\subsubsection{La Palabra Precisa}\label{la-palabra-precisa}

Aquella que \textbf{NO puede cambiarse} por otra sin alterar el
significado o matiz.

\subsubsection{Peligro de Sinónimos}\label{peligro-de-sinuxf3nimos}

``Sinónimos'' rara vez son idénticos: - Comparten zona de significado -
Mantienen diferencias de matiz - Pueden alterar registro

\subsubsection{Ejemplo}\label{ejemplo-1}

\begin{itemize}
\tightlist
\item
  ``Propósito'' (compromiso profundo)\\
  ``Fin'' (más general, menos profundo)
\end{itemize}

\subsubsection{Errores Comunes}\label{errores-comunes-4}

\begin{longtable}[]{@{}ll@{}}
\toprule\noalign{}
- Incorrecto & - Correcto \\
\midrule\noalign{}
\endhead
\bottomrule\noalign{}
\endlastfoot
Cancelar (pagar) & Pagar cuota final \\
Adolecer (carecer) & Adolecer = tener defecto \\
Cotizar (pedir) & Solicitar presupuesto \\
\end{longtable}

\subsubsection{Preferencias}\label{preferencias}

\begin{itemize}
\tightlist
\item
  \textbf{Sencilla} sobre rebuscada\\
\item
  \textbf{Corta} sobre larga\\
\item
  \textbf{Conocida} sobre rara
\end{itemize}

\subsubsection{En Ciencia}\label{en-ciencia}

\begin{itemize}
\tightlist
\item
  Tecnicismos \textbf{necesarios}
\item
  Significado \textbf{unívoco}
\item
  Evitar vaguedades: ``mejor'', ``varios'', ``proclive''
\end{itemize}

\subsubsection{Herramientas}\label{herramientas}

\begin{itemize}
\tightlist
\item
  Diccionario RAE
\item
  Diccionario de uso (María Moliner)
\item
  Diccionario ideológico (Casares)
\end{itemize}

\begin{center}\rule{0.5\linewidth}{0.5pt}\end{center}

\subsection{Conectores y Metadiscurso}\label{conectores-y-metadiscurso}

\subsubsection{Conectores: Señales de
Tránsito}\label{conectores-seuxf1ales-de-truxe1nsito}

\begin{quote}
Unidades lingüísticas que \textbf{guían inferencias} y manifiestan
relaciones lógicas.
\end{quote}

\subsubsection{Tipos Principales}\label{tipos-principales}

\paragraph{Aditivos}\label{aditivos}

además, asimismo, también, todavía más

\paragraph{Contraargumentativos}\label{contraargumentativos}

sin embargo, no obstante, por el contrario, ahora bien

\paragraph{Consecutivos}\label{consecutivos}

por lo tanto, por consiguiente, así que, entonces

\paragraph{Reformuladores}\label{reformuladores}

es decir, o sea, mejor dicho, en resumen

\paragraph{Ordenadores}\label{ordenadores}

en primer lugar, a continuación, finalmente, por último

\subsubsection{Funciones Clave}\label{funciones-clave}

\begin{itemize}
\tightlist
\item
  Garantizar ilación lógica\\
\item
  Orientar interpretación\\
\item
  Reducir ``nebulosidad''\\
\item
  Mantener coherencia
\end{itemize}

\subsubsection{Metadiscurso}\label{metadiscurso}

\textbf{Interactivo:} - ``Como se señaló anteriormente'' - ``En la
siguiente sección''

\textbf{Interaccional:} - ``Posiblemente'' - ``Considero que'' -
``Lamentablemente''

\subsubsection{Impacto}\label{impacto-1}

Construcción del \textbf{ethos} (credibilidad del autor)

\subsubsection{Evitar}\label{evitar-2}

Clichés burocráticos: - ``Es importante señalar que\ldots{}'' - ``En
otro orden de ideas\ldots{}''

\begin{center}\rule{0.5\linewidth}{0.5pt}\end{center}

\subsection{Evitar Clichés: Lenguaje
Infiltrado}\label{evitar-clichuxe9s-lenguaje-infiltrado}

\subsubsection{¿Qué son los Clichés?}\label{quuxe9-son-los-clichuxe9s}

\begin{quote}
Fórmulas que por \textbf{uso excesivo} han perdido fuerza significativa
y ``embotan la mente''.
\end{quote}

\subsubsection{Fuentes del ``Lenguaje
Infiltrado''}\label{fuentes-del-lenguaje-infiltrado}

\begin{itemize}
\tightlist
\item
  Burocracia
\item
  Medios de comunicación
\item
  Libros de texto mal escritos
\end{itemize}

\subsubsection{Ejemplos Burocráticos}\label{ejemplos-burocruxe1ticos}

\begin{itemize}
\tightlist
\item
  ``Es importante señalar que\ldots{}''\\
\item
  ``A la mayor brevedad posible''\\
\item
  ``Personalmente creemos''\\
\item
  ``El mismo/la misma'' (para evitar repetir)
\end{itemize}

\subsubsection{Clichés Literarios}\label{clichuxe9s-literarios}

\begin{itemize}
\tightlist
\item
  ``Ojos de mirar profundo''\\
\item
  ``Dientes de perla''\\
\item
  ``Venir como anillo al dedo''\\
\item
  ``Más bueno que el pan''\\
\end{itemize}

\subsubsection{Consecuencias}\label{consecuencias}

\begin{itemize}
\tightlist
\item
  Prosa con \textbf{más palabras} de lo necesario
\item
  Abuso de \textbf{voz pasiva}
\item
  Texto estático y difícil
\item
  Pérdida de \textbf{originalidad}
\end{itemize}

\subsubsection{- Alternativa}\label{alternativa}

\begin{tcolorbox}[enhanced jigsaw, arc=.35mm, rightrule=.15mm, breakable, titlerule=0mm, title=\textcolor{quarto-callout-tip-color}{\faLightbulb}\hspace{0.5em}{Cervantes sobre el Buen Español}, toprule=.15mm, left=2mm, colback=white, colbacktitle=quarto-callout-tip-color!10!white, bottomrule=.15mm, colframe=quarto-callout-tip-color-frame, toptitle=1mm, opacityback=0, opacitybacktitle=0.6, coltitle=black, bottomtitle=1mm, leftrule=.75mm]

Procurar que la oración sea ``llana'', con palabras ``honestas y bien
colocadas'', pintando la intención del autor \textbf{sin intrincar} los
conceptos.

\end{tcolorbox}

\subsubsection{Estrategia de Revisión}\label{estrategia-de-revisiuxf3n}

\begin{enumerate}
\def\labelenumi{\arabic{enumi}.}
\tightlist
\item
  Identificar expresiones vacías
\item
  Preguntar: ¿Puedo decir lo mismo con menos?
\item
  \textbf{Tachar inmisericordemente}
\item
  Asumir voz propia (uso del ``yo'')
\end{enumerate}

\subsubsection{Meta}\label{meta-2}

Texto libre de lugares comunes = más \textbf{creíble} y
\textbf{respetuoso} con el lector.

\begin{center}\rule{0.5\linewidth}{0.5pt}\end{center}

\subsection{Puntuación Correcta}\label{puntuaciuxf3n-correcta}

\subsubsection{Función}\label{funciuxf3n}

\begin{quote}
\textbf{Señales de tránsito} que mantienen flujo de información
constante y sin tropiezos.
\end{quote}

\subsubsection{Preferencia: Más Puntos, Menos
Comas}\label{preferencia-muxe1s-puntos-menos-comas}

\begin{itemize}
\tightlist
\item
  \textbf{Punto y seguido} para segmentar\\
  \textbf{Comas} solo donde necesario
\end{itemize}

\subsubsection{La Coma (,)}\label{la-coma}

\textbf{REGLA DE ORO:}\\
- NUNCA entre sujeto y verbo (salvo incisos)

\paragraph{Usos Correctos}\label{usos-correctos}

\begin{itemize}
\tightlist
\item
  Separar elementos en listas
\item
  Introducir incisos explicativos
\item
  Enmarcar marcadores (``sin embargo'')
\end{itemize}

\subsubsection{Punto y Coma (;)}\label{punto-y-coma}

\begin{itemize}
\tightlist
\item
  Pausa más larga que coma
\item
  Separar ítems complejos en listas
\item
  Vincular oraciones relacionadas
\end{itemize}

\subsubsection{Dos Puntos (:)}\label{dos-puntos}

\begin{itemize}
\tightlist
\item
  Introducir listas
\item
  Preceder ejemplos
\item
  Antes de citas directas
\item
  \begin{itemize}
  \tightlist
  \item
    NO separar verbo de objeto
  \end{itemize}
\end{itemize}

\subsubsection{Comillas (``\,``)}\label{comillas}

\begin{itemize}
\tightlist
\item
  Citas textuales
\item
  Indicar reserva/ironía
\item
  \begin{itemize}
  \tightlist
  \item
    NO para resaltar (usar sintaxis)
  \end{itemize}
\end{itemize}

\subsubsection{Impacto en Legibilidad}\label{impacto-en-legibilidad}

Puntuación deficiente: - Obliga a \textbf{releer} varias veces - Crea
\textbf{ambigüedades peligrosas} - Desmerece valor del trabajo

\subsubsection{- Revisión Final}\label{revisiuxf3n-final}

Leer en \textbf{voz alta} para verificar que pausas coinciden con
intención comunicativa.

\begin{center}\rule{0.5\linewidth}{0.5pt}\end{center}

\subsection{Legibilidad e Índice de
Nebulosidad}\label{legibilidad-e-uxedndice-de-nebulosidad}

\subsubsection{Legibilidad}\label{legibilidad}

\begin{quote}
Facilidad con la que un texto puede ser \textbf{procesado} y
\textbf{comprendido} rápidamente.
\end{quote}

\subsubsection{Dimensiones}\label{dimensiones}

\paragraph{1. Visual/Tipográfica}\label{visualtipogruxe1fica}

\begin{itemize}
\tightlist
\item
  \textbf{Letra romana} \textgreater{} Arial (remates ayudan)
\item
  \textbf{Minúsculas} \textgreater{} MAYÚSCULAS
\item
  Espacios en blanco
\end{itemize}

\paragraph{2. Léxica}\label{luxe9xica}

\begin{itemize}
\tightlist
\item
  Palabras cortas y conocidas
\item
  Evitar rebuscamiento
\end{itemize}

\paragraph{3. Estructural}\label{estructural}

\begin{itemize}
\tightlist
\item
  Oraciones: \textasciitilde18 palabras
\item
  Párrafos: \textasciitilde50 palabras
\item
  Voz activa preferida
\end{itemize}

\subsubsection{Índice de Nebulosidad (Fog
Index)}\label{uxedndice-de-nebulosidad-fog-index}

\textbf{Fórmula:}

\begin{verbatim}
A = Total palabras / Número oraciones
B = (Palabras largas × 100) / Total palabras
Índice = (A + B) × 0.4
\end{verbatim}

\textbf{Palabras largas:} \textgreater{} 3 grupos vocálicos

\subsubsection{Valores}\label{valores}

\begin{longtable}[]{@{}ll@{}}
\toprule\noalign{}
Índice & Interpretación \\
\midrule\noalign{}
\endhead
\bottomrule\noalign{}
\endlastfoot
\textbf{10} & Ideal (lectura fluida) \\
20-30 & Esfuerzo moderado \\
35+ & Esfuerzo inusual \\
84 & Textos legales (inadmisible) \\
\end{longtable}

\subsubsection{Cómo Reducir}\label{cuxf3mo-reducir}

\begin{itemize}
\tightlist
\item
  Fragmentar oraciones largas\\
\item
  Sustituir palabras complejas\\
\item
  Usar voz activa\\
\item
  Eliminar redundancia\\
\end{itemize}

\begin{center}\rule{0.5\linewidth}{0.5pt}\end{center}

\subsection{Interés Humano: El Toque
Personal}\label{interuxe9s-humano-el-toque-personal}

\subsubsection{Definición}\label{definiciuxf3n-4}

\begin{quote}
Inclusión de \textbf{personas de carne y hueso} en historias y
descripciones.
\end{quote}

\subsubsection{¿Por qué es importante?}\label{por-quuxe9-es-importante}

\begin{itemize}
\tightlist
\item
  Distingue \textbf{relato} de mero \textbf{informe}
\item
  Dispara \textbf{número de lectores}
\item
  Genera \textbf{identificación}
\item
  Facilita \textbf{comprensión} de eventos masivos
\end{itemize}

\subsubsection{Estrategias}\label{estrategias}

\begin{enumerate}
\def\labelenumi{\arabic{enumi}.}
\tightlist
\item
  \textbf{Individualizar sucesos}

  \begin{itemize}
  \tightlist
  \item
    Víctima individual \textgreater{} estadística fría
  \item
    Protagonista concreto
  \end{itemize}
\item
  \textbf{Concretar lo abstracto}

  \begin{itemize}
  \tightlist
  \item
    \begin{itemize}
    \tightlist
    \item
      ``Reinaba el nerviosismo''
    \end{itemize}
  \item
    \begin{itemize}
    \tightlist
    \item
      ``Los asistentes estaban nerviosos y\ldots{}''
    \end{itemize}
  \end{itemize}
\item
  \textbf{Humanizar la información}

  \begin{itemize}
  \tightlist
  \item
    Mostrar reacciones
  \item
    Describir personas
  \end{itemize}
\end{enumerate}

\subsubsection{Impacto}\label{impacto-2}

\begin{tcolorbox}[enhanced jigsaw, arc=.35mm, rightrule=.15mm, breakable, titlerule=0mm, title=\textcolor{quarto-callout-note-color}{\faInfo}\hspace{0.5em}{Ejemplo}, toprule=.15mm, left=2mm, colback=white, colbacktitle=quarto-callout-note-color!10!white, bottomrule=.15mm, colframe=quarto-callout-note-color-frame, toptitle=1mm, opacityback=0, opacitybacktitle=0.6, coltitle=black, bottomtitle=1mm, leftrule=.75mm]

Una tragedia con \textbf{10,000 muertos} (número) tiene menos impacto
emocional que la historia de \textbf{UNA víctima} individual.

\end{tcolorbox}

\subsubsection{En Textos Académicos}\label{en-textos-acaduxe9micos}

Aunque la objetividad es clave: - Ejemplos concretos \textgreater{}
generalidades - Casos de estudio individuales - Testimonios (en
investigación cualitativa)

\subsubsection{Balance}\label{balance}

Evitar: - Exceso de anécdota - Perder rigor científico

Mantener: - Precisión de datos - Conclusiones fundamentadas

\begin{center}\rule{0.5\linewidth}{0.5pt}\end{center}

\subsection{Uso Ético de la Inteligencia
Artificial}\label{uso-uxe9tico-de-la-inteligencia-artificial}

\subsubsection{Advertencia: ``Desinteligencia''
Artificial}\label{advertencia-desinteligencia-artificial}

\begin{quote}
La IA puede construir \textbf{``mentiras completas con verdades
parciales''}.
\end{quote}

\subsubsection{Caso Real Documentado}\label{caso-real-documentado}

ChatGPT inventó: - Título plausible de artículo - Coautores reales - DOI
falso

\textbf{Problema:} Parecía 100\% verdadero

\subsubsection{Riesgos en Formación
Académica}\label{riesgos-en-formaciuxf3n-acaduxe9mica}

\begin{itemize}
\tightlist
\item
  Desestimular lectura crítica\\
\item
  Afectar jerarquización de ideas\\
\item
  Sustituir razonamiento\\
\item
  Dependencia tecnológica\\
\item
  Erosión intelectual\\
\end{itemize}

\subsubsection{- Usos Éticos y
Eficientes}\label{usos-uxe9ticos-y-eficientes}

\textbf{Permitidos:} - Formatear referencias bibliográficas - Corrección
gramatical - Scripts de programación - Tareas tediosas técnicas

\textbf{Prohibidos:} - Generar contenido principal - Sustituir
pensamiento crítico - Inventar datos/citas - Reemplazar lectura de
fuentes

\subsubsection{Regla de Oro}\label{regla-de-oro-1}

\begin{tcolorbox}[enhanced jigsaw, arc=.35mm, rightrule=.15mm, breakable, titlerule=0mm, title=\textcolor{quarto-callout-important-color}{\faExclamation}\hspace{0.5em}{Principio Fundamental}, toprule=.15mm, left=2mm, colback=white, colbacktitle=quarto-callout-important-color!10!white, bottomrule=.15mm, colframe=quarto-callout-important-color-frame, toptitle=1mm, opacityback=0, opacitybacktitle=0.6, coltitle=black, bottomtitle=1mm, leftrule=.75mm]

Usar IA para \textbf{ahorrar tiempo} en procesos tediosos \textbf{SIN
comprometer} creatividad ni honestidad intelectual.

\end{tcolorbox}

\subsubsection{El Científico Debe\ldots{}}\label{el-cientuxedfico-debe}

Mantener capacidad de: - \textbf{Pensar} críticamente -
\textbf{Reflexionar} profundamente - \textbf{Cuestionar} activamente

\begin{center}\rule{0.5\linewidth}{0.5pt}\end{center}

\subsection{Lectura Crítica: Herramienta
Fundamental}\label{lectura-cruxedtica-herramienta-fundamental}

\subsubsection{Definición}\label{definiciuxf3n-5}

\begin{quote}
Proceso de \textbf{interpretación}, \textbf{integración de fuentes} y
\textbf{creación de información} que trasciende la decodificación.
\end{quote}

\subsubsection{Dos Niveles}\label{dos-niveles}

\paragraph{1. Lectura Extractiva}\label{lectura-extractiva}

\begin{itemize}
\tightlist
\item
  Obtención de datos aislados
\item
  Conceptos y ejemplos
\item
  Útil pero insuficiente
\end{itemize}

\paragraph{2. Lectura Reflexiva y
Crítica}\label{lectura-reflexiva-y-cruxedtica}

\begin{itemize}
\tightlist
\item
  \textbf{Sistemas de ideas} detrás de datos
\item
  Cadena de razonamientos
\item
  Relaciones lógicas entre conceptos
\item
  Posición teórica del autor
\end{itemize}

\subsubsection{Nivel Epistémico}\label{nivel-epistuxe9mico}

Leer para: - Cuestionar concepción del mundo - Transformar conocimiento
propio - Desarrollar autonomía

\subsubsection{Herramientas de Lectura
Crítica}\label{herramientas-de-lectura-cruxedtica}

\paragraph{Intertextualidad}\label{intertextualidad}

Reconocer que textos dialogan con otros

\paragraph{Campos Semánticos}\label{campos-semuxe1nticos}

Rastrear cadenas léxicas para identificar escuela teórica

\paragraph{Diccionarios
Especializados}\label{diccionarios-especializados}

\begin{itemize}
\tightlist
\item
  RAE (ortografía, existencia)
\item
  María Moliner (uso, construcción)
\item
  Casares (ideológico, sinónimos)
\end{itemize}

\subsubsection{En la Investigación}\label{en-la-investigaciuxf3n}

\textbf{Enmarcar el trabajo:} - Identificar lo novedoso - Detectar
vacíos (nicho) - Triangular evidencias

\subsubsection{Meta}\label{meta-3}

Dejar de ser \textbf{reproductor} para ser \textbf{participante activo}
en construcción de conocimiento.

\begin{center}\rule{0.5\linewidth}{0.5pt}\end{center}

\subsection{Actividad Práctica}\label{actividad-pruxe1ctica}

\subsubsection{Ejercicio 1: Análisis de
Abstract}\label{ejercicio-1-anuxe1lisis-de-abstract}

\textbf{Tiempo:} 15 minutos

\textbf{Tarea:} 1. Lee el abstract proporcionado 2. Identifica las 4
secciones IMRyD 3. Evalúa claridad, precisión, brevedad

\textbf{Criterios:} - ¿Está en un párrafo? - ¿Tiempo verbal correcto? -
¿Contiene las 4 partes? - ¿Extensión adecuada?

\subsubsection{Ejercicio 2: Reducir
Nebulosidad}\label{ejercicio-2-reducir-nebulosidad}

\textbf{Tiempo:} 10 minutos

\textbf{Tarea:} Reescribe este párrafo reduciendo su índice:

\begin{quote}
\emph{``En el transcurso del desarrollo de la investigación que fue
llevada a cabo, se procedió a la realización de un análisis exhaustivo
de los datos que fueron recolectados\ldots{}''}
\end{quote}

\subsubsection{Ejercicio 3: Identificar
Clichés}\label{ejercicio-3-identificar-clichuxe9s}

\textbf{Tiempo:} 10 minutos

\textbf{Tarea:} Encuentra y elimina clichés en:

\begin{enumerate}
\def\labelenumi{\arabic{enumi}.}
\tightlist
\item
  ``A la mayor brevedad posible''
\item
  ``Es importante señalar que''
\item
  ``Venir como anillo al dedo''
\item
  ``Personalmente creemos''
\end{enumerate}

Propón alternativas directas.

\subsubsection{Ejercicio 4: Crear Mapa de
Ideas}\label{ejercicio-4-crear-mapa-de-ideas}

\textbf{Tiempo:} 15 minutos

\textbf{Tema:} Tu proyecto de investigación

\textbf{Pasos:} 1. Escribe tema central 2. Ramifica en 3-5 ideas
principales 3. Sub-ramifica cada idea 4. Identifica conexiones

\textbf{Meta:} Base para tu Introducción

\begin{center}\rule{0.5\linewidth}{0.5pt}\end{center}

\section{Preguntas Frecuentes}\label{preguntas-frecuentes}

\subsection{❓ ¿Cuándo usar primera
persona?}\label{cuuxe1ndo-usar-primera-persona}

\textbf{R:} Depende de la revista y campo: - Métodos: Aceptable
(``Analizamos\ldots{}'') - Resultados: Preferible impersonal -
Discusión: Aceptable para interpretaciones - \textbf{Tendencia actual:}
Primera persona plural por claridad

\subsection{❓ ¿Qué extensión debe tener mi
manuscrito?}\label{quuxe9-extensiuxf3n-debe-tener-mi-manuscrito}

\textbf{R:} Varía según revista: - Artículo completo: 3,000-8,000
palabras - Nota breve: 1,500-3,000 palabras - Revisa
\textbf{instrucciones para autores}

\subsection{❓ ¿Cuántas referencias
necesito?}\label{cuuxe1ntas-referencias-necesito}

\textbf{R:} No hay número mágico: - Suficientes para respaldar
argumentos - No excesivas (denota inseguridad) - Enfoque:
\textbf{calidad} sobre cantidad

\subsection{❓ ¿Puedo usar Wikipedia como
fuente?}\label{puedo-usar-wikipedia-como-fuente}

\textbf{R:} - \textbf{NO directamente} - Úsala como punto de partida -
Verifica con fuentes primarias - Cita las fuentes originales, no
Wikipedia

\subsection{❓ ¿Cómo evito el plagio al
parafrasear?}\label{cuxf3mo-evito-el-plagio-al-parafrasear}

\textbf{R:} Técnica segura: 1. Leer y comprender la fuente 2. Cerrar el
libro/documento 3. Escribir con tus propias palabras 4. Verificar que no
sea similar 5. \textbf{SIEMPRE citar} la idea original

\subsection{❓ ¿Qué hacer si mi índice de nebulosidad es
alto?}\label{quuxe9-hacer-si-mi-uxedndice-de-nebulosidad-es-alto}

\textbf{R:} Soluciones inmediatas: - Fragmentar oraciones (\textless{}
20 palabras) - Sustituir palabras largas - Usar voz activa - Eliminar
redundancias

\begin{center}\rule{0.5\linewidth}{0.5pt}\end{center}

\section{Checklist Final para tu
Manuscrito}\label{checklist-final-para-tu-manuscrito}

\subsection{Antes de Enviar}\label{antes-de-enviar}

\subsubsection{Contenido}\label{contenido-1}

\begin{itemize}
\tightlist
\item[$\square$]
  Título atractivo y descriptivo (\textless{} 14 palabras)
\item[$\square$]
  Abstract completo con 4 secciones IMRyD
\item[$\square$]
  Introducción sigue modelo CARS
\item[$\square$]
  Métodos permiten replicación
\item[$\square$]
  Resultados objetivos (sin interpretación)
\item[$\square$]
  Discusión cierra bucle con Introducción
\item[$\square$]
  Conclusiones fundamentadas en datos
\end{itemize}

\subsubsection{Formato}\label{formato}

\begin{itemize}
\tightlist
\item[$\square$]
  Sigue estructura IMRyD
\item[$\square$]
  Referencias en formato correcto
\item[$\square$]
  Tablas y figuras numeradas
\item[$\square$]
  Pies de figura descriptivos
\item[$\square$]
  Palabras clave incluidas
\end{itemize}

\subsubsection{Estilo}\label{estilo}

\begin{itemize}
\tightlist
\item[$\square$]
  Promedio \textasciitilde18 palabras/oración
\item[$\square$]
  Párrafos 7-14 líneas
\item[$\square$]
  Voz activa predominante
\item[$\square$]
  Sin clichés ni lenguaje infiltrado
\item[$\square$]
  Conectores adecuados
\item[$\square$]
  Puntuación correcta
\item[$\square$]
  Índice nebulosidad ≤ 10
\end{itemize}

\subsubsection{Ética}\label{uxe9tica}

\begin{itemize}
\tightlist
\item[$\square$]
  Todas las fuentes citadas
\item[$\square$]
  Sin plagio (verificado)
\item[$\square$]
  Contribuciones reconocidas
\item[$\square$]
  Conflictos de interés declarados
\item[$\square$]
  Permisos obtenidos (si aplica)
\item[$\square$]
  IA usada responsablemente (si aplica)
\end{itemize}

\begin{center}\rule{0.5\linewidth}{0.5pt}\end{center}

\section{¡GRACIAS!}\label{gracias}

\subsection{Contacto y Recursos}\label{contacto-y-recursos}

\subsubsection{👨‍🏫 Edison Achalma}\label{edison-achalma}

📧 Correo institucional\\
🌐 \href{https://github.com/achalmed}{GitHub: github.com/achalmed}\\
💼 \href{https://linkedin.com/in/achalmaedison}{LinkedIn:
/in/achalmaedison}\\
🐦 \href{https://x.com/achalmaedison}{X: @achalmaedison}\\
📘 \href{https://bsky.app/profile/achalmaedison.bsky.social}{Bluesky:
achalmaedison.bsky.social}

\subsubsection{Universidad}\label{universidad}

\textbf{Corporacion Académica Universitaria}\\
Ayacucho, Perú

\subsubsection{💬 Reflexión Final}\label{reflexiuxf3n-final}

\begin{tcolorbox}[enhanced jigsaw, arc=.35mm, rightrule=.15mm, breakable, titlerule=0mm, title=\textcolor{quarto-callout-note-color}{\faInfo}\hspace{0.5em}{Recuerda}, toprule=.15mm, left=2mm, colback=white, colbacktitle=quarto-callout-note-color!10!white, bottomrule=.15mm, colframe=quarto-callout-note-color-frame, toptitle=1mm, opacityback=0, opacitybacktitle=0.6, coltitle=black, bottomtitle=1mm, leftrule=.75mm]

\begin{quote}
``La investigación científica realmente termina cuando el artículo se
\textbf{publica} en una revista sometida a revisión por pares, pues solo
entonces el hallazgo se integra al \textbf{saber científico global}.''
\end{quote}

\end{tcolorbox}

\subsubsection{🚀 ¡Éxito en tu Escritura
Científica!}\label{uxe9xito-en-tu-escritura-cientuxedfica}

\begin{center}\rule{0.5\linewidth}{0.5pt}\end{center}




\end{document}
